\documentclass[11pt,a4paper]{article}
\usepackage[margin=2.6cm]{geometry}
\usepackage[utf8]{inputenc}
\usepackage[T1]{fontenc}
\usepackage{lmodern}
\usepackage{microtype}
\usepackage{amsmath,amssymb,amsthm,mathtools}
\numberwithin{equation}{section}
\usepackage{enumitem}
\usepackage{array}
\usepackage{tabularx}
\usepackage{hyperref}
\usepackage{url}
\usepackage{tikz}
\hypersetup{colorlinks=true,linkcolor=blue,citecolor=blue,urlcolor=blue}

\theoremstyle{plain}
\newtheorem{theorem}{Theorem}[section]
\newtheorem{lemma}[theorem]{Lemma}
\newtheorem{proposition}[theorem]{Proposition}
\newtheorem{corollary}[theorem]{Corollary}
\theoremstyle{definition}
\newtheorem{definition}[theorem]{Definition}
\theoremstyle{remark}
\newtheorem{remark}[theorem]{Remark}
\newtheorem{example}[theorem]{Example}

\newcommand{\CS}{\mathrm{CS}}
\newcommand{\Cmcf}[2]{\mathcal{C}^{\mathrm{mcf}}_{#1,#2}}
\newcommand{\Clatent}[2]{\mathcal{C}^{\mathrm{mcf},\le #2}_{#1}}
\newcommand{\Callobs}[1]{\mathcal{C}^{\mathrm{mcf},\mathrm{fin}}_{#1}}
\newcommand{\D}{\mathcal{D}}
\newcommand{\out}{\mathrm{out}}
\newcommand{\tuple}[1]{\vec{#1}}
\newcommand{\blank}{\square}
\newcommand{\eps}{\varepsilon}

\providecommand{\keywords}[1]{\par\smallskip\noindent\textbf{Keywords. }#1\par}


\title{Positive-Data Learning of Multiple Context-Free Languages from Finite Observations}
\author{Takayuki Kuriyama\\\texttt{growup.kuriyama@gmail.com}}
\date{}

\newcommand{\cutpos}[1]{%
  \raisebox{-0.7ex}{%
    \tikz[baseline=(char.base)]{%
      \node[draw,circle,inner sep=0.4pt,font=\scriptsize]
        (char) {\ensuremath{#1}};%
    }%
  }%
}

\begin{document}
\setlength{\abovedisplayskip}{6pt plus 2pt minus 3pt}
\setlength{\belowdisplayskip}{6pt plus 2pt minus 3pt}
\setlength{\abovedisplayshortskip}{3pt plus 1pt minus 2pt}
\setlength{\belowdisplayshortskip}{3pt plus 1pt minus 2pt}
\maketitle


\begin{abstract}
We study positive-data learning of multiple context-free languages when the
learner is supplied with a finite compositional observation, represented by a
finite-monoid homomorphism.  For every fixed fan-out bound $f$ and explicit
observation $h$, we construct a single set-driven learner that exactly
reconstructs every language in the corresponding semantic slice from a finite
characteristic sample.  The reconstruction uses a rank search bound derived
from the observed sample and therefore requires no target-specific rule-rank
advice.  The proof combines output-type refinement, orientation propagation,
silent-child compression, and a sample-bounded witness-rank argument.  Hence
the whole slice is TxtBc-identifiable, and a conservative wrapper yields
TxtEx-identification.  The unrestricted-rank statement is qualitative; under
an additional fixed branching cap $b$, the current hypothesis is constructible
in $(1+\|K\|_+)^{O(fb)}$ time from the finite sample $K$.

The finite-observation framework is substantively broader than Yoshinaka's
multidimensional substitutable families.  On the fixed alphabet
$\Gamma=\{a,b,c\}$, the canonical symmetric boundary observation
$b_\Gamma=b_\Gamma^{1,1}$ already yields, for every $p\ge2$, an infinite
pairwise distinct family of non-context-free fan-out-two, rule-rank-one
languages in $\Cmcf{p}{b_\Gamma}\setminus\bigl(S(p)\cap p\text{-}\mathrm{MCFL}\bigr)$.  More generally, the boundary slices and finite guarded
classes satisfy a same-depth/buffered sandwich, which isolates why a genuinely
finite-state factor is needed to separate from all finite boundary guards.

We finally treat the observation as finite learning advice rather than as an
intrinsic algebraic invariant.  If it is hidden but its codomain size is
bounded by $k$, all such observations can be compiled into one effective
universal finite observation, so learnability is recovered.  Nevertheless,
these bounded latent slices already force learner-uniform
$2^{\Omega(\sqrt{k})}$ positive locking information along an explicit family.
If no bound on the hidden finite observation is supplied, TxtBc-identifiability
fails.  A Rambow--Satta calibration separately exhibits unbounded minimum rule
rank already over one fixed three-letter alphabet and shows that the finite
observation size required by this construction still grows with the rank
parameter.
\end{abstract}

\keywords{grammatical inference, multiple context-free grammars, positive data, finite observations, multidimensional substitutability, TxtBc learning, TxtEx learning, rule rank}

\section{Introduction}
\label{sec:introduction}

Positive-data learning from text originates with Gold and has developed into a
large theory of grammatical inference from incomplete evidence
\cite{Gold1967,Angluin1980,Angluin1982Reversible}.  A particularly successful
route is distributional: substrings or tuples that occur in sufficiently
similar contexts are treated as interchangeable.  For context-free languages
this viewpoint underlies substitutable-language learning, and Yoshinaka
extended it to multiple context-free languages (MCFLs) by using
multidimensional contexts \cite{ClarkEyraud2007,Yoshinaka2008,Yoshinaka2009,Yoshinaka2011}.
The latter work is the closest predecessor of the present paper.

A basic difficulty for unrestricted MCFG learning is that fan-out does not by
itself bound the number of nonterminals that may occur on the right-hand side
of a rule.  Yoshinaka's positive-data learner therefore fixes both a dimension
(or fan-out) parameter $p$ and a rule-rank parameter $q$.  This raises a
natural qualitative question: can one identify a semantically defined MCFG
family without supplying a target-specific rule-rank bound to the learner?
The issue is not whether each target has finite rule rank---every finite MCFG
presentation does---but whether the learner must know such a bound in advance.
At the qualitative level, however, removal of an a priori rank parameter is
not by itself a separation from Yoshinaka's framework.  His fixed-rank
conjecture construction is uniformly sound in the rank cap $q$ on positive
samples from a $p$D-substitutable target---that is, increasing $q$ does not
introduce words outside the target language---and its characteristic-sample
completeness applies whenever the chosen rank cap is at least the rank of a
target presentation \cite{Yoshinaka2011}; see also the explicit restatement in
\cite[Proposition~1]{CaseYoshinakaZeugmann2013}.  Consequently the cap can be allowed to grow with the
sample; Section~\ref{sec:comparison} records this adaptive-rank baseline
explicitly.  The contribution below is therefore not the bare learnability of
$\bigcup_q SL(p,q)$, but a canonical reconstruction theorem for the broader
finite-observation semantic slices introduced next.

The finite-observation idea has a fan-out-one predecessor in the author's
context-free setting \cite{kuriyamaCFG}.%
\footnote{This result originated in the author's master's thesis at the
Graduate University for Advanced Studies (SOKENDAI) and underwent internal
academic review; see
\url{https://makotokanazawa.ws.hosei.ac.jp/students-j.html}.  A revised version
is \cite{kuriyamaCFG}.  The results of the present paper do not depend on this
earlier result.}
The multidimensional extension, however, faces a structural obstruction absent
for ordinary CFGs.  At fan-out one,
standard CFG binarization removes right-hand-side rank as an independent
parameter.  For MCFGs of fan-out at least two, by contrast, minimum rule rank
forms a genuine hierarchy \cite{RambowSatta1999}.  Thus the
present result is not a tuple-valued restatement of the CFG theorem: the proof
must track tuple orientations and induced child orientations through
compositions, compress genuinely silent children, and recover the necessary
unbounded search rank from positive sample support.

We answer this question by supplying a different finite interface.  A
\emph{finite observation} is an explicit finite-monoid homomorphism
$h:\Sigma^*\to M$.  Tuple occurrences are also equipped with their
left-to-right orientation in the surrounding sentence.  Roughly, two tuples
may be merged when they have the same orientation, equal componentwise
$h$-values, and share an accepting context of that orientation.  The resulting
semantic class at fan-out at most $f$ is denoted $\Cmcf{f}{h}$.  The map $h$
is neither a parse tree nor a membership oracle: it is a fixed finite-state
map evaluated uniformly on strings.  It can arise, for example, from a known
regular envelope or from the alphabet-dependent boundary observation used in
our comparison with Yoshinaka.

We regard $h$ as class-level finite side information, not as a canonical
invariant that must be recovered separately for each target.  In particular,
the same observation is shared by every language in one slice, and our main
strict-separation example below uses a single observation for an infinite
family of target languages.  Sections~\ref{sec:fixed-observation} and
\ref{sec:comparison} then make this modeling choice testable: a bound on the
size of a hidden observation can be compiled into universal finite advice,
whereas removing every such bound destroys TxtBc-identifiability.

The characteristic-sample theorem below is not obtained by effectively
filtering the natural space of MCFG presentations to those satisfying the
semantic promise.  Appendix~\ref{app:decision-complexity} shows that this
recognition problem is $\Pi^0_1$-complete, already for context-free
presentations with the trivial observation.  This statement concerns the
natural grammar presentation of the target class; it does not assert the
nonexistence of every alternative effective indexing.

Our main result gives exact reconstruction on every fixed observation slice.
Starting from an arbitrary target presentation only on the proof side, we
normalize it without increasing fan-out or rule rank, refine productive states
by their componentwise $h$-types, and propagate tuple orientations through the
normalized rules.  Genuinely silent child occurrences are then compressed.
Every remaining child in a successful rule exposure contributes positive
terminal support; consequently a rule witnessed inside a sample word $w$ has
rank at most $|w|$.  The learner can therefore enumerate all necessary
composition witnesses using a rank cap derived from the current finite sample,
not from the unknown target presentation.  A finite characteristic sample
exposes every normalized target rule, while the semantic substitution rules
remain sound.  Thus one set-driven learner exactly reconstructs every
$L\in\Cmcf{f}{h}$ and yields both TxtBc and TxtEx identification.  If an
external branching cap $b$ is fixed, the corresponding current hypothesis is
constructible in $(1+\|K\|_+)^{O(fb)}$ time.  This latter statement concerns
hypothesis construction from a supplied finite sample; it is not a polynomial
bound on the total length of characteristic data.

The observation can also be viewed as finite learning advice, and this gives a
sharp qualitative boundary.  Suppose the actual morphism is hidden but its
codomain size is known to be at most $k$.  Because only finitely many monoid
structures and letter maps of size at most $k$ exist, they can be compiled into
one effective universal observation $U_{\Sigma,k}$; refinement monotonicity
then reduces the bounded latent problem to the supplied-observation theorem.
This recovery is not quantitatively free: even for finite regular targets,
any TxtBc learner for the whole bounded slice has locking sequences whose
positive content is $2^{\Omega(\sqrt{k})}$ along an explicit family.  If the
observation is completely latent and no size bound is given, the union of all
finite-observation slices is not TxtBc-identifiable from positive data.

The comparison with Yoshinaka calibrates both the scope and the significance
of the result.  On the fixed alphabet $\Gamma=\{a,b,c\}$, the canonical
symmetric boundary observation $b_\Gamma=b_\Gamma^{1,1}$ already satisfies,
for every $p\ge2$,
\[
  S(p)\cap p\text{-}\mathrm{MCFL}
  \subsetneq
  \Cmcf{p}{b_\Gamma},
\]
and the difference contains an infinite pairwise distinct family
$H_t=\{a^{tn}b^nc^n:n\ge0\}$ of non-context-free fan-out-two languages of
minimum rule rank one.  Thus the substantive separation already occurs in the
simplest canonical alphabet-only slice, rather than only after adding an
auxiliary regular-control factor.  At the same time, the comparison with
finite $(k,\ell)$ boundary guards has a sharp limitation: the same-depth
boundary slice is contained in the corresponding guarded semantic class,
whereas one additional boundary symbol gives the converse embedding.  A fixed
regular-control observation is therefore retained below for the stronger,
different purpose of separating from every finite boundary guard.

We separately note that the qualitative learnability of $\bigcup_q SL(p,q)$
already follows from an adaptive-rank use of Yoshinaka's own conjecture
construction; it is not claimed as a contribution here.  A tagged
Rambow--Satta construction instead serves as a structural calibration.  Its
exact minimum-rank hierarchy can be compressed to the fixed three-letter
alphabet $\{a,b,c\}$ without collapsing rank, so the unboundedness is not an
artifact of alphabet growth.  Nevertheless a shared-context type-packing
argument survives this compression: every finite observation witnessing the
parameter-$r$ fixed-alphabet language still satisfies
$r\le |\operatorname{im}h|^2$.  Thus the construction does not establish
unbounded rank inside one fixed observation slice, and the lower bound is not
asserted as a general rank--observation tradeoff.  Finally, the output-type
construction yields a rank-preserving decomposition of each fixed-observation
target into finitely many congruential components, providing a direct bridge
to the syntactic and distributional approaches of Clark and Yoshinaka.

The contributions can therefore be summarized as follows.
\begin{enumerate}[leftmargin=*]
\item \textbf{Exact reconstruction on fixed-observation slices.}
For fixed $f$ and explicit $h$, every language in $\Cmcf{f}{h}$ has a finite
characteristic sample for one effective canonical learner; the slice is both
TxtBc- and TxtEx-identifiable.  The required search rank is derived from the
sample rather than supplied as target-specific advice.  Under a fixed external
branching cap $b$, the current capped hypothesis is constructible in
$(1+\|K\|_+)^{O(fb)}$ time.
\item \textbf{A substantive separation in the canonical boundary slice.}
On the fixed alphabet $\Gamma=\{a,b,c\}$, the single canonical observation
$b_\Gamma=b_\Gamma^{1,1}$ works for every $p\ge2$ and yields a separation
whose difference contains infinitely many pairwise distinct non-context-free
fan-out-two languages of minimum rule rank one.  The corresponding finite
boundary slices also satisfy a same-depth/buffered sandwich with guarded
multidimensional substitutability, clarifying exactly where genuinely
finite-state control becomes stronger than finite boundary information.
\item \textbf{Observation-advice boundary.}
A bound $k$ on a hidden finite observation suffices via a universal refinement;
bounded latent slices already impose a learner-uniform
$2^{\Omega(\sqrt{k})}$ locking lower bound, whereas the unbounded latent union
is not TxtBc-identifiable.
\item \textbf{Structural calibration.}
An alphabet-only boundary observation already contains Yoshinaka's whole
semantic family.  A Rambow--Satta minimum-rank hierarchy remains exact after
compression to one fixed three-letter alphabet, but every finite observation
witnessing its parameter-$r$ member has image size at least $\sqrt r$; hence
that construction still does not place unbounded rank inside one fixed
observation slice.  Each fixed-observation target also admits a rank-preserving
finite decomposition into congruential components.
\end{enumerate}

The paper is organized as follows.  Section~\ref{sec:prelim} fixes the learning
model, MCFG notation, oriented tuple distributions, and finite-observation
semantic slices.  Section~\ref{sec:running-examples} gives two guiding
boundaries.  Section~\ref{sec:learner} constructs the canonical learner and
proves exact reconstruction, and Section~\ref{sec:poly-time} isolates the
fixed-rank construction bound.  Section~\ref{sec:fixed-observation} studies
bounded and unbounded latent observations from the learning-theoretic side.
Section~\ref{sec:comparison} compares the framework with Yoshinaka and related
distributional approaches.  Appendix~\ref{app:rank-hierarchy} gives the
technical rank-hierarchy calibration, and Appendix~\ref{app:decision-complexity}
records the decision complexity of the semantic promise.

\section{Learning Model and Finite-Observation Semantics}
\label{sec:prelim}

\subsection{Positive-Data Learning}

\begin{definition}[Text and TxtBc/TxtEx identification]
To handle the empty language, fix a pause symbol
\(\#\notin\Sigma\).%
\footnote{Allowing the pause symbol makes
\(\#,\#,\ldots\) a text for the empty language.}
A \emph{text} for \(L\subseteq\Sigma^*\) is an infinite sequence over
\(L\cup\{\#\}\) in which every element of \(L\) appears at least once.
Pause symbols may occur arbitrarily, and \(\#\) is not an end-of-text marker.

The \emph{content} of a finite prefix is the finite set
\(K\subseteq L\) of positive examples from \(\Sigma^*\) that occur in that
prefix.  Thus neither the positions nor the number of pause symbols, nor repeated
occurrences of the same positive example, affect the content.
A learner is a computable function mapping each finite prefix to a hypothesis
grammar.

A learner \emph{TxtBc-identifies} a language \(L\) from positive data if,
for every text for \(L\), there exists \(n_0\) such that every hypothesis
produced from stage \(n_0\) onward has language \(L\)
\cite{CaseLynes1982,CaseSmith1983}.  Thus eventual semantic correctness of
every hypothesis is required, but syntactic stabilization of the hypothesis
grammar is not.

A learner \emph{TxtEx-identifies} \(L\) from positive data if, for every text
for \(L\), there are \(n_0\) and a single hypothesis grammar \(H\) such that
all stages from \(n_0\) onward output exactly \(H\), with \(L(H)=L\).  A
learner TxtBc- or TxtEx-identifies a class when it does so for every language
in the class.  We prove the TxtBc result with a set-driven learner and recover
TxtEx convergence by a conservative wrapper in
Corollary~\ref{cor:gold-explanatory}.  Unqualified ``identification in the
limit'' is avoided below when the distinction matters.
\end{definition}

\begin{definition}[Set-driven learner and characteristic sample]
A learner \(\mathcal A\) is \emph{set-driven} if its output hypothesis is
independent of presentation order and multiplicity and depends only on the
finite set \(K\) of positive examples observed so far.  Write
\(\mathcal A(K)\) for the hypothesis on observation set \(K\).

For a target language \(L\), a finite set \(S\subseteq L\) is a
\emph{characteristic sample for \(\mathcal A\)} if, for every finite set
\(K\), \(S\subseteq K\subseteq L\) implies \(L(\mathcal A(K))=L\).
The case \(S=\emptyset\) is allowed.

The canonical learner constructed below is set-driven.  Characteristic samples
are therefore the natural quantitative certificates for the positive results
and for set-driven lower bounds.  Arbitrary learners may additionally exploit presentation order, multiplicity,
and pauses.  A learner-uniform lower bound against such learners will use
semantic locking sequences only when needed in Section~\ref{sec:fixed-observation}.
The presentation-relative characteristic samples constructed below are
concrete characteristic samples for the canonical learner.  The target
presentation is used only to establish their existence; it is not supplied as
input to the learner.
\end{definition}

The semantic locking notion below, and the corresponding existence argument,
are standard in behaviorally correct language learning; compare the TxtBc
framework of Case and Lynes \cite{CaseLynes1982} and subsequent treatments.
We include the short proof for self-containedness.  The new use made in this
paper is the deletion-family application in
Section~\ref{subsec:uniform-observation-bound}, which turns locking content
into a learner-uniform quantitative lower bound for the bounded-observation
MCFG classes studied here.

\begin{definition}[Positive sample size]
For finite \(K\subseteq\Sigma^*\), define
\(\|K\|_+:=\sum_{w\in K}\max(1,|w|)\).%
\footnote{The factor \(\max(1,|w|)\) prevents an observed empty word from
making a degenerate zero contribution to sample size.}
For a tuple \(\tuple u=(u_1,\ldots,u_d)\), write
\(\|\tuple u\|_1:=\sum_{i=1}^d |u_i|\)
for its total component length.
\end{definition}

The general MCFG theorem below allows \(\eps\)-generating rank-zero rules.

\begin{remark}[Slicewise polynomiality]
\label{rem:slicewise-poly}
Unless a theorem makes additional parameter dependence explicit, polynomial
bounds are understood slicewise in the fixed external parameters.  In
particular, Theorem~\ref{thm:poly} fixes the fan-out bound, a branching cap,
and the supplied finite observation; no polynomial bound uniform in those
parameters is claimed.
\end{remark}

\subsection{Finite Observations}

\begin{definition}[Explicit finite homomorphism]
\label{def-monoid-morphism}
A monoid homomorphism \(h\colon\Sigma^*\to M\) into a finite monoid \(M\) is
\emph{explicit} if the multiplication table of \(M\) and the values \(h(a)\) for
\(a\in\Sigma\) are given.  For a tuple \(\tuple w=(w_1,\ldots,w_d)\), write
\(h^{(d)}(\tuple w):=(h(w_1),\ldots,h(w_d))\in M^d\).
\end{definition}

In this paper, finite-monoid homomorphisms are used as finite-state auxiliary
observations for positive-data learning of MCFLs.  For the classical
correspondence between finite-monoid homomorphisms and finite-automaton
recognition of regular languages, see Pin~\cite{Pin1986}.

A finite observation \(h\) may arise, for example, when a DFA recognizing a
regular language \(R\supseteq L\) is available in advance.  The transition
morphism of this DFA yields an explicit homomorphism \(h:\Sigma^*\to M\): for
each word \(w\in\Sigma^*\), the element \(h(w)\) records the state
transformation induced by reading \(w\).

Thus \(h(u)=h(v)\) means that \(u\) and \(v\), when read from any DFA state,
lead to the same resulting state.  The learner does not infer \(R\) itself and
does not use \(R\) as a membership oracle for \(L\).  It uses only this finite
state information from the DFA as auxiliary information controlling when
tuples observed in positive data may be identified.


\subsection{Multiple Context-Free Grammars}

We use the standard tuple-generating definition of multiple context-free
grammars \cite{SekiEtAl1991}; see also \cite{Kallmeyer2010} for a general
overview of MCFGs and related mildly context-sensitive formalisms.
The explicit general-rank notation below is used throughout the learner, the
proofs, and the running examples.

\begin{definition}[Standard linear MCFG presentation]
\label{def:standard-mcfg}
A \emph{multiple context-free grammar} is a tuple
\(G=(V,\Sigma,P,S,\mu)\), where \(V\) is a finite set of nonterminals,
\(\Sigma\) is a finite terminal alphabet, \(S\in V\) has \(\mu(S)=1\), and
\(\mu(A)\ge1\) is the fan-out of each \(A\in V\).  A rule of rank \(r\ge0\)
has the form
\[
  \rho:\quad
  A\to(\alpha_1,\ldots,\alpha_e)(B_1,\ldots,B_r),
\]
where \(e=\mu(A)\), \(d_j=\mu(B_j)\), and, with
\[
  X_j:=\{x_j^1,\ldots,x_j^{d_j}\},
  \qquad
  \Gamma_\rho:=\Sigma\uplus X_1\uplus\cdots\uplus X_r,
\]
each \(\alpha_i\in\Gamma_\rho^*\) and every variable in
\(X_1\uplus\cdots\uplus X_r\) occurs at most once in the complete tuple
\((\alpha_1,\ldots,\alpha_e)\).  The rule is \emph{nondeleting} if every such
variable occurs exactly once.

For compatible child tuples
\(\tuple u_j=(u_{j,1},\ldots,u_{j,d_j})\), simultaneous substitution of
\(u_{j,k}\) for \(x_j^k\) defines
\[
  \rho(\tuple u_1,\ldots,\tuple u_r)\in(\Sigma^*)^e.
\]
When \(r=0\), the rule is a constant rule and derives the tuple
\((\alpha_1,\ldots,\alpha_e)\).  The tuple languages \(L_A(G)\) form the least family, ordered componentwise by
inclusion, that is closed under all rule applications, and
\[
  L(G):=\{w\in\Sigma^*:(w)\in L_S(G)\}.
\]
An \(f\)-MCFG has \(\mu(A)\le f\) for every nonterminal, and an \(f\)-MCFL is
a language generated by an \(f\)-MCFG.  A grammar is \(b\)-branching if every
rule has rank at most \(b\).
\end{definition}

For integers \(f,q\ge1\), write
\[
  L(f,q):=\{L(G):G\text{ is an MCFG of fan-out at most }f
  \text{ and rule rank at most }q\}.
\]
Thus \(f\text{-}\mathrm{MCFL}=\bigcup_{q<\infty}L(f,q)\), because every
finite MCFG presentation has finite rule rank.
For \(L\in f\text{-}\mathrm{MCFL}\), define its \emph{minimum rule rank at
fan-out \(f\)} by
\[
  \operatorname{rr}_f(L)
  :=\min\{q\ge1:L\in L(f,q)\}.
\]
This is a language-level minimum over all fan-out-\(f\) presentations, not
the rank of a particular supplied grammar.


When two nonterminals have the same fan-out \(d\), we use
\[
  A\to B
\]
as shorthand for the unary identity-template rule
\(A\to(x_1^1,\ldots,x_1^d)(B)\).%
\footnote{This abbreviation denotes only the identity correspondence between
components.  For example, at fan-out two, \(A\to B\) means
\(A\to(x_1^1,x_1^2)(B)\); it does not denote a component-swapping rule such as
\(A\to(x_1^2,x_1^1)(B)\).}
Likewise, \(A\to\tuple c\) denotes a rank-zero
constant rule.  These abbreviations are used for typed start links and for the
semantic unary rules of the learner.

The standard formalism allows rank zero, nonidentity unary rules, arbitrary
finite rule rank, and deleting rules.  A linear nondeleting rule
\(\rho:A\to(\alpha_1,\ldots,\alpha_e)(B_1,\ldots,B_r)\) is
\emph{nonpermuting} if, for each child \(B_j\), the variables
\(x_j^1,\ldots,x_j^{\mu(B_j)}\) occur in this order when the parent template
components \(\alpha_1,\ldots,\alpha_e\) are scanned from left to right.
We use a standard rank-preserving normalization to separate the semantic
theorem from presentation artifacts.  Kanazawa explicitly records that every
$r$-ary branching $m$-MCFG has an equivalent $r$-ary branching $m$-MCFG whose
rules are all nondeleting and nonpermuting
\cite[Section~2]{Kanazawa2019Ogden}.  Thus fan-out and the branching cap may be
preserved simultaneously.  For the quantitative bookkeeping used below, we
realize this normalization in two transparent steps.  First, the standard
deletion-elimination construction for MCFGs
\cite[Lemma~2.2]{SekiEtAl1991} yields an equivalent nondeleting presentation
without increasing fan-out or branching rank.  In outline, one decorates a
nonterminal by the subset of output components retained by a surrounding
derivation and projects each rule to the variables that actually survive.  A
completely unused child occurrence is removed rather than copied; no new child
occurrence is introduced, so rule rank cannot increase, and every retained
tuple has arity at most that of the original nonterminal, so fan-out cannot
increase.  Starting from a nondeleting presentation,
Lemma~\ref{lem:permutation-decoration} below gives a self-contained
permutation-decoration construction producing an equivalent nonpermuting
presentation while again preserving fan-out and maximum rule rank.  Second,
rule rank is a genuine independent parameter at bounded fan-out: the hierarchy
by independent parallelism does not collapse when synchronized parallelism is
bounded \cite{RambowSatta1999}.  At the grammar-transformation level, LCFRS
binarization may require increased fan-out \cite{GomezRodriguezEtAl2009}.
For fan-out two, Sagot and Satta give an optimal procedure for reducing the
right-hand-side rank of a given LCFRS production while respecting the fan-out
two constraint \cite{SagotSatta2010}; this is a presentation-level
factorization problem and does not by itself determine the minimum rank among
all grammars generating the same language.%
\footnote{For intuition at the grammar level, consider fan-out-two
nonterminals \(A,B,C,D,E\) and the rank-four rule
\(A\to(x_1^1x_2^1x_3^1x_4^1,\,
 x_3^2x_1^2x_4^2x_2^2)(B,C,D,E)\).
For every pair of children, the corresponding variables are adjacent in one
parent component but separated by variables of other children in the other.
Thus any first binary combination must retain at least three pieces if it is
to support the later interleaving.  This illustrates why a particular rule
need not admit a fan-out-two strong binarization; the language-level
non-collapse is supplied by Rambow and Satta \cite{RambowSatta1999}.}
Thus the arbitrary-rank witness construction below is not replaceable by a
harmless binary normal-form step when the fan-out parameter \(f\) is fixed.

\begin{definition}[Productivity, reachability, and reduction]
\label{def:productive-reachable-reduced}
Let \(G=(V,\Sigma,P,S,\mu)\) be a standard MCFG.

A nonterminal \(A\in V\) is \emph{productive} if
\(L_A(G)\ne\emptyset\).

The \emph{nonterminal dependency graph} of \(G\) has vertex set \(V\) and,
for every rule
\(A\to(\alpha_1,\ldots,\alpha_e)(B_1,\ldots,B_k)\)
and every \(1\le j\le k\), an edge \(A\to B_j\).
A nonterminal \(A\) is \emph{reachable} if there is a directed path from
\(S\) to \(A\) in this graph.  Paths of length zero are allowed, so the start
symbol \(S\) is always reachable.

The grammar \(G\) is \emph{reduced} if every nonterminal is both reachable
and productive.
\end{definition}

\paragraph{Rule rank versus fan-out.}
The \emph{rank} of a production is the number of nonterminal children on its
right-hand side, whereas the fan-out \(\mu(A)\) is the arity of tuples derived
from \(A\).  These notions are independent.  For example, if
\(\mu(A)=\mu(B)=3\), the unary rule
\(A\to(x_1^1a,x_1^2,bx_1^3)(B)\) has rank one while \(A\) has fan-out three.
Conversely, if \(\mu(C)=\mu(D)=\mu(E)=\mu(F)=1\), then
\(C\to(x_1^1x_2^1x_3^1)(D,E,F)\) has parent fan-out one but rule rank three.
Fan-out controls the arity of tuple objects, while rule rank controls how many
child occurrences must be combined simultaneously in a composition witness.

For a rule
\(\rho:A\to(\alpha_1,\ldots,\alpha_e)(B_1,\ldots,B_r)\), we call
\((\alpha_1,\ldots,\alpha_e)\) its \emph{template tuple}.  All examples below
use the standard MCFG notation of Definition~\ref{def:standard-mcfg}; no
separate binary grammar formalism is assumed.

\subsection{Oriented Tuple Contexts and Semantic Slices}

\begin{definition}[Sentence context and orientation]
\label{def:sentence-context}
A \emph{sentence context of arity \(d\)} is a word of the form
\[
  E=u_0\blank_{\sigma(1)}u_1\cdots\blank_{\sigma(d)}u_d,
\]
where \(\sigma\) is a permutation of \(\{1,\ldots,d\}\), each
\(u_i\in\Sigma^*\), and each named hole \(\blank_i\) occurs exactly once.
Its \emph{orientation} is
\(\operatorname{ori}(E):=\sigma\).
For \(\tuple w=(w_1,\ldots,w_d)\), write \(E[\tuple w]\) for the string
obtained by substituting \(w_i\) for \(\blank_i\).
For a permutation \(\sigma\in S_d\), let \(\mathcal E_{d,\sigma}\) denote
the set of arity-\(d\) sentence contexts of orientation \(\sigma\).
\end{definition}

\begin{definition}[Induced child orientation]
\label{def:induced-orientation}
Let
\(\rho:A\to(\alpha_1,\ldots,\alpha_e)(B_1,\ldots,B_r)\)
be a linear nondeleting rule template, and let
\(\varpi\in S_e\) be an orientation of its output tuple.  For a child \(B_j\)
of fan-out \(d_j\), scan the concatenated template components
\[
  \alpha_{\varpi(1)}\alpha_{\varpi(2)}\cdots\alpha_{\varpi(e)}
\]
from left to right and record only the variables
\(x_j^1,\ldots,x_j^{d_j}\) belonging to child \(j\).  Linearity and
nondeletion make their superscripts appear exactly once each, hence in a
unique permutation of \(\{1,\ldots,d_j\}\).  Denote this permutation by
\[
  \operatorname{ind}_j(\rho,\varpi)\in S_{d_j}.
\]
It is the orientation induced on child \(j\) when the parent tuple is exposed
with orientation \(\varpi\).
\end{definition}

For a nondeleting rule \(\rho\), the nonpermuting condition introduced above
is equivalently
\[
  \operatorname{ind}_j(\rho,\mathrm{id}_{\mu(A)})
  =\mathrm{id}_{\mu(B_j)}
  \qquad(1\le j\le r).
\]

\begin{lemma}[Permutation-decoration normalization]
\label{lem:permutation-decoration}
Let \(G\) be a finite linear nondeleting \(f\)-MCFG.  There is an equivalent
nonpermuting \(f\)-MCFG \(G^{\mathrm{np}}\) with the same maximum rule rank
such that each original nonterminal has at most \(f!\) decorated copies and
each original rule has at most \(f!\) decorated copies.  Hence
\[
  |G^{\mathrm{np}}|=O_f(|G|).
\]
\end{lemma}

\begin{proof}
For every nonterminal \(A\) of fan-out \(d\), introduce copies
\(A^\tau\) indexed by \(\tau\in S_d\); the intended tuple language of
\(A^\tau\) is obtained from \(L_A(G)\) by permuting components by \(\tau\).
For an original rule and a chosen permutation of its parent components, scan
the permuted parent template.  For each child, the scan determines a unique
permutation of that child's component variables.  Use the corresponding
decorated child symbol and rename those variables into natural order.  The
resulting rule is nonpermuting by construction.  Starting from the identity
copy of the original start symbol gives a derivation-by-derivation
correspondence with \(G\).

There are at most \(d!\le f!\) copies of any symbol, and choosing the parent
permutation determines all child decorations, so the same bound holds per
rule.  No rule gains or loses children; hence rule rank is preserved.
Removing unreachable decorated copies preserves equivalence and cannot
increase fan-out or rule rank.
\end{proof}


\begin{definition}[Tuple occurrence in a sample word]
\label{def:tuple-occurrence}
Let \(K\subseteq\Sigma^*\) and let \(w=a_1\cdots a_n\in K\).
Number the cut positions of \(w\) as
\(\cutpos{0}a_1\cutpos{1}a_2\cutpos{2}\cdots a_n\cutpos{n}\).
For \(0\le\ell\le r\le n\), write
\(w[\ell:r]:=a_{\ell+1}\cdots a_r\); in particular,
\(w[\ell:\ell]=\eps\).

Let \(d\ge1\).  An \emph{arity-\(d\) tuple occurrence} in \(w\) is data
\(\mathfrak{o}=(\sigma,(\ell_i,r_i)_{i=1}^d)\), where \(\sigma\in S_d\) and
\[
  0\le\ell_{\sigma(1)}\le r_{\sigma(1)}
  \le\cdots\le\ell_{\sigma(d)}\le r_{\sigma(d)}\le n.
\]
Empty intervals \(\ell_i=r_i\) are allowed, and when several empty intervals
occupy the same cut position, \(\sigma\) also records their local order.

For such an occurrence, put
\(\tuple x=(w[\ell_1:r_1],\ldots,w[\ell_d:r_d])\).
Replacing the specified intervals by their corresponding named holes yields an
arity-\(d\) context \(E\) with \(\operatorname{ori}(E)=\sigma\) and
\(E[\tuple x]=w\).  We identify the occurrence with the oriented pair
\((E,\tuple x)\).

For example, let
\(w=\cutpos{0}a\cutpos{1}b\cutpos{2}c\cutpos{3}d\cutpos{4}\), let \(d=2\),
and take \((\ell_1,r_1)=(2,4)\) and \((\ell_2,r_2)=(0,1)\).
The second component occurs first from the left, so \(\sigma=(2,1)\); the
tuple is \((cd,a)\), and the sentence context is \(\blank_2b\blank_1\).
\end{definition}

Template components in MCFG rules may be empty, which is why
Definition~\ref{def:tuple-occurrence} allows empty intervals.  For example,
for \(w=ab\), if both components of \((\eps,\eps)\) are observed at the
middle cut, the two local orders \(\mathrm{id}_2\) and \((2,1)\) produce the distinct
contexts \(a\blank_1\blank_2b\) and \(a\blank_2\blank_1b\).  Orientation is
therefore genuine finite occurrence information even when tuple values agree.

\begin{definition}[Sector tuple distribution]
\label{def:tuple-distribution}
For \(L\subseteq\Sigma^*\), \(\tuple x\in(\Sigma^*)^d\), and
\(\sigma\in S_d\), define the \emph{orientation-sector distribution}
\[
  \D_{L,\sigma}^{(d)}(\tuple x)
  :=
  \{E\in\mathcal E_{d,\sigma}:E[\tuple x]\in L\}.
\]
The full named-context distribution is the disjoint sector union
\[
  \D_L^{(d)}(\tuple x)
  :=\bigcup_{\sigma\in S_d}\D_{L,\sigma}^{(d)}(\tuple x).
\]
When the arity is clear, superscripts are omitted.
\end{definition}

\begin{definition}[Orientation-aware \((f,h)\)-tuple substitutability]
\label{def:tuple-substitutable}
Let \(f\ge1\) and let \(h\colon\Sigma^*\to M\) be an explicit finite monoid
homomorphism.  A language \(L\subseteq\Sigma^*\) is
\emph{orientation-aware \((f,h)\)-tuple-substitutable} if, for every
\(1\le d\le f\), every \(\sigma\in S_d\), and all
\(\tuple x,\tuple y\in(\Sigma^*)^d\),
\[
 h^{(d)}(\tuple x)=h^{(d)}(\tuple y),\qquad
 \D_{L,\sigma}^{(d)}(\tuple x)\cap
 \D_{L,\sigma}^{(d)}(\tuple y)\ne\emptyset
\]
imply
\[
  \D_{L,\sigma}^{(d)}(\tuple x)
  =\D_{L,\sigma}^{(d)}(\tuple y).
\]
We usually shorten this to \emph{\((f,h)\)-tuple-substitutable}.  Thus the
orientation is part of the finite structural state and substitutability is
required inside, but not across, permutation sectors.  Recall that
\(h^{(d)}(\tuple w)=(h(w_1),\ldots,h(w_d))\in M^d\).
\end{definition}

\begin{lemma}[Canonicalization of orientation sectors]
\label{lem:sector-canonicalization}
Let \(d\ge1\) and \(\sigma\in S_d\).  For
\[
  \tuple x=(x_1,\ldots,x_d),
\]
put
\[
  P_\sigma(\tuple x):=(x_{\sigma(1)},\ldots,x_{\sigma(d)}).
\]
For a context
\[
  E=u_0\blank_{\sigma(1)}u_1\cdots\blank_{\sigma(d)}u_d
  \in\mathcal E_{d,\sigma},
\]
define
\[
  C_\sigma(E):=u_0\blank_1u_1\cdots\blank_du_d
  \in\mathcal E_{d,\mathrm{id}_d}.
\]
Then \(C_\sigma\) is a bijection
\[
  \mathcal E_{d,\sigma}\longrightarrow\mathcal E_{d,\mathrm{id}_d}
\]
and
\[
  C_\sigma(E)[P_\sigma(\tuple x)]=E[\tuple x].
\]
Consequently,
\[
  C_\sigma\!\left(\D_{L,\sigma}^{(d)}(\tuple x)\right)
  =
  \D_{L,\mathrm{id}_d}^{(d)}\!\left(P_\sigma(\tuple x)\right).
\]
\end{lemma}

\begin{proof}
The map \(C_\sigma\) merely renames holes according to their left-to-right
order.  It is therefore bijective, with inverse obtained by renaming
\(\blank_i\) back to \(\blank_{\sigma(i)}\).  The filling identity follows
directly from the definitions, and restricting the same bijection to
accepting contexts gives the distribution identity.
\end{proof}

\begin{corollary}[Identity-sector formulation]
\label{cor:identity-sector-equivalence}
Definition~\ref{def:tuple-substitutable} is equivalent to requiring its
substitutability implication only in the identity sector
\(\sigma=\mathrm{id}_d\), for every \(1\le d\le f\) and all tuples of
arity \(d\).
\end{corollary}

\begin{proof}
One direction is immediate.  Conversely, fix \(\sigma\in S_d\) and tuples
\(\tuple x,\tuple y\) satisfying the premises of
Definition~\ref{def:tuple-substitutable} in sector \(\sigma\).  Applying the
same permutation to both tuples preserves equality of componentwise
observation values.  Lemma~\ref{lem:sector-canonicalization} transports their
shared accepting context to the identity sector.  The identity-sector
assumption therefore gives equality of the two identity-sector distributions,
and transport back by \(C_\sigma^{-1}\) gives
\[
  \D_{L,\sigma}^{(d)}(\tuple x)
  =
  \D_{L,\sigma}^{(d)}(\tuple y).
\]
\end{proof}

\begin{remark}[Sectorwise versus full-context substitutability]
\label{rem:full-context-stronger}
Lemma~\ref{lem:sector-canonicalization} shows that the individual orientation
sectors are equivalent up to a permutation of tuple coordinates.  Thus
quantifying over all sectors in Definition~\ref{def:tuple-substitutable} is a
coordinate-covariant way of stating the condition, rather than an additional
semantic restriction beyond the identity sector.

This should be distinguished from the genuinely stronger requirement obtained
by replacing each sector distribution \(\D_{L,\sigma}^{(d)}\) by the full
distribution \(\D_L^{(d)}\).  The latter couples contexts belonging to
different orientation sectors and therefore cannot be reduced to a coordinate
renaming.  Proposition~\ref{prop:orientation-strict} shows that this stronger
implication is strict.  We retain explicit sector labels later only as
proof-side bookkeeping for coordinate order in composition witnesses; by
Corollary~\ref{cor:identity-sector-equivalence}, they are neither additional
learning advice nor an additional semantic restriction.
\end{remark}

\begin{proposition}[Sectorwise substitutability is strictly weaker than full-context substitutability]
\label{prop:orientation-strict}
There are a finite alphabet \(\Sigma\), a finite regular language
\(L\subseteq\Sigma^*\), and an explicit finite observation
\(h:\Sigma^*\to M\) such that \(L\) is orientation-aware
\((2,h)\)-tuple-substitutable but fails the stronger full-context condition.
Hence the implication in Remark~\ref{rem:full-context-stronger} is strict.
\end{proposition}

\begin{proof}
Let \(\Sigma=\{a,b,c,d\}\) and
\[
  L:=\{ac,bd,ca\}.
\]
Put \(\Gamma=\{A,C\}\), map \(a,b\) to \(A\) and \(c,d\) to \(C\), and let
\(T_2(\Gamma)=\Gamma^{\le2}\cup\{\top\}\) be the Rees quotient monoid that
records a word of length at most two exactly and sends every longer product to
the absorbing element \(\top\).  Composing the letter map with the quotient
morphism gives a finite observation \(h:\Sigma^*\to T_2(\Gamma)\).

We first verify the orientation-aware condition.  Suppose tuples
\(\tuple x,\tuple y\) of arity \(d\le2\) have equal componentwise
\(h\)-type and share an accepting context \(E\) in one fixed orientation.
Because every word of \(L\) has length two, no component occurring in such a
filling can have \(h\)-value \(\top\); corresponding components therefore
have the same length and the same \(A/C\)-pattern.  If the total tuple length
is two, every accepting context in the fixed orientation has empty terminal
material, hence that sector contains only the corresponding pure-hole
context.  If the total tuple length is one, a distinct equal-type replacement
could only exchange \(a\) with \(b\), or \(c\) with \(d\).  No one-letter
context can complete both members of either pair to words of \(L\): the
possible completions of \(a\) and \(b\) are different, and likewise for
\(c\) and \(d\).  Thus the shared-context premise forces the relevant
one-letter components to agree.  If the total tuple length is zero, the tuples
are identical.  In every case the two sector distributions are equal, so
\(L\) is orientation-aware \((2,h)\)-tuple-substitutable.

For failure of the full-context condition, take
\[
  \tuple x=(a,c),\qquad \tuple y=(b,d).
\]
Their componentwise \(h\)-types agree.  In the identity sector,
\(E=\blank_1\blank_2\) accepts both tuples because
\(E[\tuple x]=ac\in L\) and \(E[\tuple y]=bd\in L\).  In the reverse sector,
however, \(F=\blank_2\blank_1\) satisfies
\[
  F[\tuple x]=ca\in L,
  \qquad
  F[\tuple y]=db\notin L.
\]
Hence the full named-context distributions differ although they intersect,
which is exactly a failure of the stronger condition.  Thus strictness already
occurs for a finite regular language.
\end{proof}

For \(d=1\) there is only one orientation, and sentence contexts are ordinary
two-sided contexts \(u\blank_1v\).  Hence
Definition~\ref{def:tuple-substitutable} specializes exactly to the fixed-\(h\)
two-sided substitutability condition for CFLs \cite{kuriyamaCFG}.

\begin{definition}[Observation--distribution equivalence]
\label{def:observation-distribution-equivalence}
Let $L\subseteq\Sigma^*$, let $h:\Sigma^*\to M$ be a finite observation,
let $d\ge1$, and let $\sigma\in S_d$.  For tuples
$\tuple x,\tuple y\in(\Sigma^*)^d$, write
\[
  \tuple x\equiv_{L,h,\sigma}^d\tuple y
\]
if and only if
\[
  h^{(d)}(\tuple x)=h^{(d)}(\tuple y)
  \qquad\text{and}\qquad
  \D_{L,\sigma}^{(d)}(\tuple x)=\D_{L,\sigma}^{(d)}(\tuple y).
\]
For fixed $L,h,d,\sigma$ this is an equivalence relation.  When the arity is
clear, the superscript $d$ may be omitted.
\end{definition}

\begin{definition}[The target class]
\label{def:class-cmcf}
Fix \(f\ge1\) and an explicit finite monoid homomorphism \(h\).  The class
\(\Cmcf{f}{h}\) consists of all languages \(L\subseteq\Sigma^*\) such that
\[
  L\in f\text{-}\mathrm{MCFL}
  \qquad\text{and}\qquad
  L\text{ is orientation-aware }(f,h)\text{-tuple-substitutable}.
\]
\end{definition}

\begin{proposition}[$h$-recognizable languages lie in the fixed-observation slice]
\label{prop:h-recognizable-in-slice}
Let $h:\Sigma^*\to M$ be a finite monoid homomorphism and let $P\subseteq M$.
Then, for every $f\ge1$,
\[
  h^{-1}(P)\in\Cmcf{f}{h}.
\]
Thus every fixed-observation slice contains the full Boolean algebra of
languages recognized by its observation morphism.
\end{proposition}

\begin{proof}
Put $L=h^{-1}(P)$.  Since $h$ is finite, $L$ is regular and hence an
$f$-MCFL for every $f\ge1$.  Fix an arity $d\le f$, an orientation
$\sigma\in S_d$, and tuples $\tuple x,\tuple y$ with
$h^{(d)}(\tuple x)=h^{(d)}(\tuple y)$.  In any fixed named sentence context
of orientation $\sigma$, the homomorphism property gives the same total
$h$-value after filling with $\tuple x$ or $\tuple y$, because corresponding
components have equal $h$-values.  Hence the two fillings are either both in
$h^{-1}(P)$ or both outside it.  Their sector distributions are therefore
equal; the shared-context premise is not needed.  Thus $L$ is
orientation-aware $(f,h)$-tuple-substitutable.
\end{proof}

For fixed \(f\) and \(h\), we call \(\Cmcf{f}{h}\) the
\emph{fixed-observation slice determined by \(h\)}.  The observation morphism
is fixed externally and shared by all targets in the class; orientation is
not additional advice, because it is read directly from each concrete tuple
occurrence.

\begin{remark}[Semantic nature of the target class]
\label{rem:semantic-class}
Membership in \(\Cmcf{f}{h}\) is a semantic promise imposed on the target
language.  The learner is defined on every finite positive sample, but it does
not receive an MCFG presentation of the target language and does not test this
promise from positive data.  This is not merely a modeling convenience: Appendix~\ref{app:decision-complexity}
shows that recognition of the promise from an MCFG presentation and an explicit
observation is $\Pi^0_1$-complete, already in the context-free case with the
trivial observation.  Related decision problems for distributional and
congruential properties have been studied in the context-free setting
\cite{ClarkEyraud2007,Clark2017,KanazawaKappe2019}.
\end{remark}

\begin{definition}[Refinement of finite observation morphisms]
\label{def:h-refinement}
Let \(h\colon\Sigma^*\to M\) and \(h'\colon\Sigma^*\to M'\) be explicit finite
homomorphisms.  We say that \(h'\) \emph{refines} \(h\), and write
\(h\preceq h'\), if there is a monoid homomorphism
\(\pi\colon M'\to M\) such that \(h=\pi\circ h'\).
\end{definition}

\begin{proposition}[Monotonicity under refinement of the observation morphism]
\label{prop:h-refinement-monotonicity}
If \(h\preceq h'\) and \(L\) is \((f,h)\)-tuple-substitutable, then \(L\) is
\((f,h')\)-tuple-substitutable.  Consequently, for every pair of explicit
finite homomorphisms with \(h\preceq h'\),
\(\Cmcf{f}{h}\subseteq \Cmcf{f}{h'}\).
\end{proposition}

\begin{proof}
Let \(\pi\colon M'\to M\) satisfy \(h=\pi\circ h'\).  Fix an arity
\(d\le f\) and an orientation \(\sigma\in S_d\).  Suppose that
\(h'^{(d)}(\tuple x)=h'^{(d)}(\tuple y)\) and that the two tuples share an
accepting context of orientation \(\sigma\).  Applying \(\pi\) componentwise
gives \(h^{(d)}(\tuple x)=h^{(d)}(\tuple y)\).  The
\((f,h)\)-tuple-substitutability of \(L\) therefore gives
\(\D_{L,\sigma}^{(d)}(\tuple x)=\D_{L,\sigma}^{(d)}(\tuple y)\), which is
precisely the \((f,h')\)-condition in this sector.  Membership in
\(f\text{-}\mathrm{MCFL}\) is unchanged, so the class inclusion follows.
\end{proof}

\begin{lemma}[Shared-context substitutability]
\label{lem:shared-context}
Let \(L\) be \((f,h)\)-tuple-substitutable.  If \(d\le f\),
\(\sigma\in S_d\), \(h^{(d)}(\tuple x)=h^{(d)}(\tuple y)\), and a concrete
context \(E\) of orientation \(\sigma\) satisfies
\(E[\tuple x],E[\tuple y]\in L\), then
\[
  \tuple x\equiv_{L,h,\sigma}^d\tuple y.
\]
\end{lemma}

\begin{proof}
The context \(E\) belongs to both sector distributions.  Definition~\ref{def:tuple-substitutable}
gives equality of those distributions, and the type equality is assumed.
\end{proof}



\section{Guiding Examples}
\label{sec:running-examples}

The examples below serve two purposes.  The three-block language shows how a
coarse regular observation can coexist with genuinely multiple-context-free
numerical synchronization.  The copy language marks the opposite boundary:
being a $2$-MCFL alone does not guarantee the existence of any finite
observation witnessing the semantic promise.

\begin{proposition}[Regular control and an abelian balance constraint]
\label{prop:regular-abelian-balance}
Let \(R\subseteq\Sigma^*\) be a regular language, represented by a finite
monoid \(M\), a homomorphism \(h_R\colon\Sigma^*\to M\), and a subset
\(P\subseteq M\) such that \(R=h_R^{-1}(P)\).
Let \(\mathbb A\) be an abelian group written additively,
let \(\varphi\colon\Sigma^*\to\mathbb A\) be a monoid homomorphism, and let
\(H\le\mathbb A\) be a subgroup.
Then
\(L:=R\cap\varphi^{-1}(H)\)
is \((f,h_R)\)-tuple-substitutable for every \(f\ge1\).
\end{proposition}

\begin{proof}
It is enough to prove the stronger full-context condition of
Remark~\ref{rem:full-context-stronger}.  Fix \(1\le d\le f\), and let
\(\tuple{x}=(x_1,\ldots,x_d)\) and
\(\tuple{y}=(y_1,\ldots,y_d)\) satisfy
\(h_R^{(d)}(\tuple{x})=h_R^{(d)}(\tuple{y})\) and
\(\D_L^{(d)}(\tuple{x})\cap\D_L^{(d)}(\tuple{y})\ne\emptyset\).
It suffices to show
\(\D_L^{(d)}(\tuple{x})=\D_L^{(d)}(\tuple{y})\).

Fix
\(E\in\D_L^{(d)}(\tuple{x})\cap\D_L^{(d)}(\tuple{y})\).
Then \(E[\tuple{x}],E[\tuple{y}]\in L\), and in particular
\(\varphi(E[\tuple{x}]),\varphi(E[\tuple{y}])\in H\).

First consider the regular control.
Let
\(F=u_0\blank_{\sigma(1)}u_1\cdots
\blank_{\sigma(d)}u_d\)
be an arbitrary arity-\(d\) sentence context.
Since
\(h_R^{(d)}(\tuple{x})=h_R^{(d)}(\tuple{y})\),
we have \(h_R(x_i)=h_R(y_i)\) for every \(i\), and hence
\(h_R(F[\tuple{x}])=h_R(F[\tuple{y}])\).
Because \(R=h_R^{-1}(P)\),
\(F[\tuple{x}]\in R\) if and only if \(F[\tuple{y}]\in R\).

Next consider the abelian balance constraint.
Put
\(\Delta:=\sum_{i=1}^d(\varphi(y_i)-\varphi(x_i))\in\mathbb A\).
Since \(\mathbb A\) is abelian, the order of the holes is irrelevant, and for every
arity-\(d\) sentence context \(F\),
\(\varphi(F[\tuple{y}])-\varphi(F[\tuple{x}])=\Delta\).
The fixed terminal parts of \(F\) cancel, leaving only the sum of the
componentwise differences.

Taking \(F=E\) gives
\(\Delta=\varphi(E[\tuple{y}])-\varphi(E[\tuple{x}])\).
Both terms lie in \(H\), so \(\Delta\in H\).
Hence for every \(F\),
\(\varphi(F[\tuple{y}])=\varphi(F[\tuple{x}])+\Delta\), and therefore
\(F[\tuple{x}]\in\varphi^{-1}(H)\) if and only if
\(F[\tuple{y}]\in\varphi^{-1}(H)\).

Combining this with \(L=R\cap\varphi^{-1}(H)\), we obtain, for every
arity-\(d\) sentence context \(F\),
\(F[\tuple{x}]\in L\) if and only if \(F[\tuple{y}]\in L\).
Thus
\(\D_L^{(d)}(\tuple{x})=\D_L^{(d)}(\tuple{y})\).

Since \(d\le f\) was arbitrary, \(L\) is
\((f,h_R)\)-tuple-substitutable.
\end{proof}


For the positive block envelope
\(P_m:=a_1^+a_2^+\cdots a_m^+\),
let \(D_m^+\) be the complete DFA with state set
\(\{0,1,\ldots,m,\mathsf{sink}\}\),
initial state \(0\), and unique accepting state \(m\).

For each \(1\le i<m\), use the transitions
\[
  i\xrightarrow{a_i}i,
  \qquad
  i\xrightarrow{a_{i+1}}i+1,
\]
together with
\(0\xrightarrow{a_1}1\) and
\(m\xrightarrow{a_m}m\).
All transitions not listed above go to \(\mathsf{sink}\), and
\(\mathsf{sink}\) loops on every letter.

Let
\(h_m^+:\Sigma_m^*\to M_m^+\)
be the transition morphism of \(D_m^+\).
Then \(D_m^+\) recognizes exactly \(P_m\).


\begin{example}[The three-block agreement language]
\label{ex:l3-language}
Let
\(L_3:=\{a^n b^n c^n\mid n\ge1\}\) and
\(P_3:=a^+b^+c^+\), and put \(h_{P_3}:=h_3^+\).
This morphism records the coarse block zone of a component but not the equality
of the three block lengths.
\end{example}

\begin{proposition}[A compact fan-out-two presentation for \(L_3\)]
\label{prop:l3-running-example}
The language \(L_3\) has a reduced linear nondeleting MCFG presentation of
fan-out two and rule rank two.
\end{proposition}

\begin{proof}
Let \(V=\{S,T,A,A_a,A_b,A_c\}\), let \(\mu(A)=2\), and let all other
nonterminals have fan-out \(1\).
Use
\[
  S\to T,\quad
  A_a\to(a),\quad
  A_b\to(b),\quad
  A_c\to(c),\quad
  A\to(x_1^1,x_2^1)(A_a,A_b),
\]
\[
  A\to(ax_1^1,bx_1^2x_2^1)(A,A_c),\qquad
  T\to(x_1^1x_1^2x_2^1)(A,A_c).
\]

A direct induction gives
\[
  L_A=\{(a^n,b^nc^{n-1})\mid n\ge1\}.
\]
Hence the top rule
\(T\to(x_1^1x_1^2x_2^1)(A,A_c)\),
together with \(S\to T\), generates exactly \(L_3\).

Finally, every nonterminal is reachable and productive, and every rule is
linear and nondeleting.
Therefore the grammar is a reduced linear nondeleting MCFG of fan-out two and rule rank two.
\end{proof}

\begin{corollary}[The three-block example belongs to the fixed-observation class]
\label{cor:l3-in-fixed-observation-class}
For every \(f\ge2\), \(L_3\in\Cmcf{f}{h_{P_3}}\).
\end{corollary}

\begin{proof}
The presentation is given by Proposition~\ref{prop:l3-running-example}.
For the regular control \(P_3=a^+b^+c^+\), define
\(\varphi(w):=(|w|_a-|w|_b,\ |w|_b-|w|_c)\in\mathbb Z^2\).
Then \(L_3=P_3\cap\varphi^{-1}(\{(0,0)\})\), so
Proposition~\ref{prop:regular-abelian-balance} gives
\((f,h_{P_3})\)-tuple substitutability.
\end{proof}




\begin{proposition}[The copy language has no finite witnessing observation]
\label{prop:copy-infinite-observation}
Let \(\Sigma:=\{a,b\}\) and define
\(\mathrm{COPY}:=\{ww:w\in\Sigma^+\}\).
Then \(\mathrm{COPY}\in2\text{-}\mathrm{MCFL}\), but for every finite
monoid homomorphism \(h:\Sigma^*\to M\), the language \(\mathrm{COPY}\) is
not \((2,h)\)-tuple-substitutable.
\end{proposition}

\begin{proof}
We first show that \(\mathrm{COPY}\in2\text{-}\mathrm{MCFL}\).
Use a nonterminal \(A\) of fan-out \(2\) and a start symbol \(S\) of fan-out
\(1\).  Give \(A\) the rules
\(A\to(a,a)\), \(A\to(b,b)\),
\(A\to(a x_1^1,a x_1^2)(A)\), and
\(A\to(b x_1^1,b x_1^2)(A)\), and use the start rule
\(S\to(x_1^1x_1^2)(A)\).
Then \(L_A=\{(w,w):w\in\Sigma^+\}\).  Indeed, the two base rules derive
\((a,a)\) and \((b,b)\), while each recursive rule maps an already derived
\((w,w)\) to either \((aw,aw)\) or \((bw,bw)\), and no other form can be
derived.  Thus the start rule generates exactly
\(\{ww:w\in\Sigma^+\}\), proving that
\(\mathrm{COPY}\in2\text{-}\mathrm{MCFL}\).

Now fix an arbitrary finite monoid \(M\) and an arbitrary monoid homomorphism
\(h\colon\Sigma^*\to M\).  We show that \(\mathrm{COPY}\) is not
\((2,h)\)-tuple-substitutable.  Choose \(n\ge1\) with \(2^n>|M|\).
Since there are exactly \(2^n\) words of length \(n\), the pigeonhole
principle yields distinct \(u,v\in\Sigma^n\) with \(h(u)=h(v)\).

Let \(\tuple x:=(u,u)\) and \(\tuple y:=(v,v)\).  Then
\(h^{(2)}(\tuple x)=h^{(2)}(\tuple y)\).  Moreover, for the arity-two named
context \(E:=\blank_1\blank_2\), we have
\(E[\tuple x]=uu\in\mathrm{COPY}\) and
\(E[\tuple y]=vv\in\mathrm{COPY}\).  Both \(E\) and the context used below have orientation \(\mathrm{id}_2\), and
\[
E\in\D_{\mathrm{COPY},\mathrm{id}_2}^{(2)}(\tuple x)\cap
\D_{\mathrm{COPY},\mathrm{id}_2}^{(2)}(\tuple y).
\]

The two distributions are nevertheless different already inside this same
orientation sector.  Let \(F:=\blank_1\blank_2uu\).  Then
\(F[\tuple x]=uuuu=(uu)(uu)\in\mathrm{COPY}\).  On the other hand,
\(F[\tuple y]=vvuu\).  Suppose for contradiction that
\(vvuu\in\mathrm{COPY}\).  Then \(vvuu=zz\) for some \(z\in\Sigma^+\).
Since \(|vvuu|=4n\), we have \(|z|=2n\).  Thus the first and second halves
of \(vvuu\), each of length \(2n\), must be equal, so \(vv=uu\).
Because \(|u|=|v|=n\), comparing the first \(n\) symbols gives \(u=v\), a
contradiction.  Therefore \(F[\tuple y]\notin\mathrm{COPY}\).
Consequently
\(\D_{\mathrm{COPY},\mathrm{id}_2}^{(2)}(\tuple x)\neq
\D_{\mathrm{COPY},\mathrm{id}_2}^{(2)}(\tuple y)\).

Thus \(\tuple x\) and \(\tuple y\) have the same componentwise \(h\)-type
and share an accepting context in one orientation sector, but have different
distributions in that same sector.
By Definition~\ref{def:tuple-substitutable}, \(\mathrm{COPY}\) is not
\((2,h)\)-tuple-substitutable.  Since the finite-monoid homomorphism \(h\)
was arbitrary, no finite observation witnesses \((2,h)\)-tuple
substitutability for \(\mathrm{COPY}\).
\end{proof}

\section{Canonical Reconstruction from Positive Data}
\label{sec:learner}

\paragraph{Proof architecture.}
The reconstruction proof has four dependencies.  First, output-type refinement
and orientation propagation put every productive proof state into one fixed
componentwise \(h\)-type and one exposed orientation.  Second, silent-child
compression removes children that contribute no terminal material without
changing fan-out or the target language.  Third, every remaining child in a
successful derivation has positive support, so a rule witnessed inside a sample
word has rank at most the positive sample size.  Finally, the semantic unary
rules identify exactly those observed tuples whose common-context evidence and
finite types justify substitution.  Soundness uses these ingredients
child-by-child; completeness chooses a finite characteristic sample exposing
every normalized target rule.  The dependency chain is summarized by
\[
\begin{aligned}
  \text{typing/orientation}
  &\Longrightarrow \text{silent compression}
  \Longrightarrow \text{sample-bounded rank}\\
  &\Longrightarrow \text{soundness + finite rule exposure}
  \Longrightarrow \text{exact reconstruction}.
\end{aligned}
\]


The learner uses only the fixed fan-out bound \(f\), the supplied observation
\(h\), and the finite positive sample available at the current stage.  A target
MCFG is used only on the proof side to establish the existence of a finite
characteristic sample.  The target presentation is never supplied to the
learner.

For the qualitative theorem, let \(G\) be any finite standard MCFG of fan-out
at most \(f\) whose rules are nondeleting and nonpermuting.  In the final
exact-reconstruction theorem such a presentation is obtained, without
changing fan-out or branching rank, from an arbitrary standard MCFG
presentation by the normalization recalled in Section~\ref{sec:prelim}.
By Corollary~\ref{cor:identity-sector-equivalence}, one could equivalently
canonicalize every observed tuple into left-to-right component order and
formulate the semantic promise using only the identity sector.  We retain
explicit orientation annotations in the learner because they preserve the
original component names of sample occurrences and make induced child
orientations under arbitrary sample-derived templates transparent.  Thus
orientation is bookkeeping for composition, not additional target advice or
a stronger semantic assumption.

The proof architecture is worth making explicit.  Starting from an arbitrary
standard \(f\)-MCFG presentation of the target language, we first pass on the
proof side to an equivalent nondeleting, nonpermuting presentation without
increasing fan-out or branching rank.  We then refine nonterminals by their
finite observation output types and eliminate genuinely silent child
occurrences by compression.  Because the proof-side rules are nonpermuting,
identity orientation propagates from the root to every child occurrence; this
lets the characteristic-sample and completeness arguments use one canonical
identity-sector exposure per proof-side state.  The resulting positive-rank
rules are exposed directly, without binarization.  Finally, although their
ranks are not bounded uniformly over the target class, every rule needed for
reconstruction has a positive sample exposure whose length bounds that
particular rank.  Thus the general learner can enumerate all required
witnesses using only its sample-dependent rank cap.

\subsection{Output-Type Refinement and Exposing Contexts}
\label{sec:refinement}

\begin{definition}[Template output type]
\label{def:output-map}
Let
\(\rho:A\to(\alpha_1,\ldots,\alpha_e)(B_1,\ldots,B_r)\)
be a nondeleting rule, put \(e=\mu(A)\) and
\(d_j:=\mu(B_j)\), and choose
\(\mathbf q_j\in M^{d_j}\) for each \(1\le j\le r\).

On the template alphabet
\(
  \Gamma_\rho
  :=\Sigma\uplus\biguplus_{j=1}^r
  \{x_j^1,\ldots,x_j^{d_j}\}
\),
define
\(\theta_{\mathbf q_1,\ldots,\mathbf q_r}\colon\Gamma_\rho\to M\)
by
\(\theta_{\mathbf q_1,\ldots,\mathbf q_r}(a):=h(a)\)
for \(a\in\Sigma\), and
\(\theta_{\mathbf q_1,\ldots,\mathbf q_r}(x_j^k):=(\mathbf q_j)_k\).

Let
\(
  \widehat\theta_{\mathbf q_1,\ldots,\mathbf q_r}
  \colon\Gamma_\rho^*\to M
\)
be its unique extension to a monoid homomorphism.
Define
\[
  \out_\rho(\mathbf q_1,\ldots,\mathbf q_r)
  :=
  \bigl(
    \widehat\theta_{\mathbf q_1,\ldots,\mathbf q_r}(\alpha_1),
    \ldots,
    \widehat\theta_{\mathbf q_1,\ldots,\mathbf q_r}(\alpha_e)
  \bigr)
  \in M^e.
\]
For rank zero, \(\out_\rho\) is simply
\((h(\alpha_1),\ldots,h(\alpha_e))\), the componentwise \(h\)-type of the
constant template.
\end{definition}

\begin{definition}[Output-type refinement]
\label{def:output-refinement}
Given a finite nondeleting MCFG \(G=(V,\Sigma,P,S,\mu)\) and a finite
observation \(h\), the \emph{full output-type refinement} \(G^h\) has a fresh
start symbol \(\widetilde S\) and typed nonterminals
\(A_{\mathbf p}\) for every \(A\in V\) and
\(\mathbf p\in M^{\mu(A)}\).
Set
\(\mu(\widetilde S):=1\) and
\(\mu(A_{\mathbf p}):=\mu(A)\).

For each \(m\in M\), add the start rule
\(\widetilde S\to S_{(m)}\).
For each rank-zero rule \(A\to\tuple c\), add
\(A_{h^{(\mu(A))}(\tuple c)}\to\tuple c\).

For every positive-rank rule
\(\rho:A\to(\alpha_1,\ldots,\alpha_e)(B_1,\ldots,B_r)\)
and every choice
\(\mathbf q_j\in M^{\mu(B_j)}\) (\(1\le j\le r\)), add
\[
  A_{\out_\rho(\mathbf q_1,\ldots,\mathbf q_r)}
  \to
  \rho(B_{1,\mathbf q_1},\ldots,B_{r,\mathbf q_r}).
\]
Here \(B_{j,\mathbf q_j}\) denotes the typed copy of \(B_j\) indexed by
\(\mathbf q_j\).

Finally, the \emph{trimmed refinement} \(\widetilde G_0\) is obtained by
retaining only the typed nonterminals and typed rules that actually occur in
some successful derivation from \(\widetilde S\).
\end{definition}

\begin{proposition}[Output-type invariants]
\label{prop:output-invariants}
For \(G^h\) and \(\widetilde G_0\):
\begin{enumerate}[label=(\roman*),leftmargin=*]
\item if \(A_{\mathbf p}\Rightarrow^*\tuple u\), then
\(h^{(\mu(A))}(\tuple u)=\mathbf p\), and erasing the type indices gives a
valid \(G\)-derivation of \(\tuple u\);
\item every fixed \(G\)-derivation tree lifts uniquely to \(G^h\) by labeling
each node with the componentwise \(h\)-type of the tuple derived at that node;
\item \(L(\widetilde G_0)=L(G^h)=L(G)\).
\end{enumerate}
\end{proposition}

\begin{proof}
At a rule application, homomorphic evaluation of the rule template gives
\[
  h^{(\mu(A))}
  \bigl(\rho(\tuple u_1,\ldots,\tuple u_r)\bigr)
  =
  \out_\rho
  \bigl(h^{(\mu(B_1))}(\tuple u_1),\ldots,
        h^{(\mu(B_r))}(\tuple u_r)\bigr).
\]
The rank-zero case is immediate.  Induction on derivation height proves~(i),
and the same identity gives the unique lifting in~(ii).  Erasing types gives
one language inclusion, lifting gives the other, and trimming removes only
objects absent from successful derivations.
\end{proof}

\begin{lemma}[Output typing preserves nonpermuting rules]
\label{lem:refinement-nonpermuting}
If the proof-side grammar \(G\) is nonpermuting, then every non-start rule of
\(G^h\), and hence every such rule of the trimmed refinement
\(\widetilde G_0\), is nonpermuting.
\end{lemma}

\begin{proof}
Output typing changes only nonterminal labels and leaves each rule template
unchanged.  The relative order of every child's component variables is
therefore unchanged.
\end{proof}

For a surviving typed nonterminal \(X=A_{\mathbf p}\), put
\(L_X:=\{\tuple u:X\Rightarrow_{\widetilde G_0}^*\tuple u\}\).
It is nonempty by trimming.

\begin{definition}[Proof-side identity exposing context]
\label{def:identity-exposing-context}
Let \(X=A_{\mathbf p}\) survive in \(\widetilde G_0\), with arity \(d\).
An \emph{identity exposing context} for \(X\) is a sentence context
\(E\in\mathcal E_{d,\mathrm{id}_d}\) such that
\[
  E[\tuple u]\in L(G)
  \qquad\text{for every }\tuple u\in L_X.
\]
\end{definition}

\begin{lemma}[Identity exposures exist in the normalized proof grammar]
\label{lem:identity-exposing-contexts}
Assume the proof-side grammar \(G\) is nondeleting and nonpermuting.  Every
surviving typed nonterminal \(X\) has an identity exposing context.  If
\(\widetilde S\to X\) is a typed start rule, one may use \(\blank_1\).
\end{lemma}

\begin{proof}
Choose a successful typed derivation containing an occurrence of \(X\), delete
the subtree at that occurrence, and replace its output components by named
holes.  At the root the single component has identity orientation.  By
Lemma~\ref{lem:refinement-nonpermuting}, every typed rule is nonpermuting, so
Definition~\ref{def:induced-orientation} gives identity orientation to each
child whenever its parent is identity-oriented.  Induction down the selected
root-to-\(X\) path therefore shows that the deleted occurrence of \(X\) is
exposed in identity orientation.  Replacing the deleted subtree by any other
\(X\)-derivation yields another successful derivation, so the resulting
context exposes all of \(L_X\).
\end{proof}

For every surviving typed nonterminal \(X\), fix one identity exposing context
and write it \(\chi(X)\).  Whenever \(X\) is the child of a surviving typed
start rule \(\widetilde S\to X\), choose \(\chi(X)=\blank_1\).  For every
remaining \(X\), choose any identity exposing context obtained from a
successful derivation as in the preceding lemma.  These choices are
proof-side existence witnesses; they are neither supplied to nor computed by
the learner.

\begin{lemma}[Identity orientation propagates through normalized rules]
\label{lem:identity-orientation-propagation}
Let \(R:X\to\rho(Y_1,\ldots,Y_r)\) be a surviving typed rule.  If \(G\) is
nonpermuting, then
\[
  \operatorname{ind}_j(\rho,\mathrm{id}_{\mu(X)})
  =\mathrm{id}_{\mu(Y_j)}
  \qquad(1\le j\le r).
\]
The same statement holds for retained children of a silently compressed rule.
\end{lemma}

\begin{proof}
The first assertion is exactly the nonpermuting condition, preserved by
output-type refinement.  Silent-child compression removes variables of
children that derive only empty strings and leaves the relative order of all
retained child variables unchanged.
\end{proof}

\subsection{Silent-Child Compression and Finite Rule Exposure}

\begin{definition}[Silent typed nonterminal and positive anchor]
\label{def:silent-typed}
A surviving typed nonterminal \(X\) of arity \(d\) is \emph{silent} if
\(L_X=\{(\eps,\ldots,\eps)\}\).

If \(X\) is not silent, put
\(L_X^+:=\{\tuple u\in L_X:\|\tuple u\|_1>0\}\), and choose
\(\omega(X)\in L_X^+\) of minimum total length, breaking ties by a fixed
shortlex order on tuple encodings.  If \(X\) is silent, put
\(\omega(X):=(\eps,\ldots,\eps)\).
\end{definition}

\begin{lemma}[Identity-sector equivalence of the empty tuple and the positive anchor]
\label{lem:mixed-empty-positive}
Let \(L:=L(G)\) be \((f,h)\)-tuple-substitutable, and let \(X\) be a
surviving typed nonterminal of arity \(d\).  If \(X\) is not silent and
\((\eps,\ldots,\eps)\in L_X\), then
\[
  \omega(X)\equiv_{L,h,\mathrm{id}_d}^d(\eps,\ldots,\eps).
\]
\end{lemma}

\begin{proof}
Both tuples lie in \(L_X\), hence have the same componentwise output type by
Proposition~\ref{prop:output-invariants}.  The identity exposing context
\(\chi(X)\) accepts both, so Lemma~\ref{lem:shared-context} applies in the
identity sector.
\end{proof}

\begin{definition}[Silent-child compression]
\label{def:silent-compression}
Keep every fresh start link \(\widetilde S\to S_{(m)}\) unchanged.  For every
other surviving typed rule \(R:X\to\rho(Y_1,\ldots,Y_r)\), delete from the
body every child occurrence that is silent and delete its variables from the
template.  Retain the remaining non-silent children in their original order
and call the resulting rule \(R^+:X\to\rho^+(Y_{j_1},\ldots,Y_{j_s})\).
If all children are silent, \(s=0\); in that case \(\rho^+\) is the
rank-zero constant tuple obtained from the original template after deleting
all variables of silent children.  Remove duplicate resulting rules and call
the compressed grammar \(\widetilde G_0^+\).
\end{definition}

\begin{lemma}[Silent-child compression preserves all typed tuple languages]
\label{lem:silent-compression}
The grammar \(\widetilde G_0^+\) is a finite linear nondeleting MCFG with the
same nonterminals and fan-outs as \(\widetilde G_0\).  If the proof-side grammar
\(G\) is nonpermuting, then \(\widetilde G_0^+\) is nonpermuting as well.  Moreover,
\(L_X(\widetilde G_0^+)=L_X(\widetilde G_0)\) for every surviving typed
nonterminal \(X\).  In particular \(L(\widetilde G_0^+)=L(G)\).  If
\(X=A_{\mathbf p}\) is the left-hand side of a compressed rule, every tuple
produced by that rule has componentwise \(h\)-type \(\mathbf p\).
\end{lemma}

\begin{proof}
A deleted silent child derives only the all-empty tuple, whose componentwise
\(h\)-type is the identity tuple.  Removing all variables of such a child
therefore changes neither any ground parent tuple nor its output type.  Notice
that a non-silent typed state that can also derive the all-empty tuple is
\emph{not} deleted: it is retained with its positive anchor
\(\omega(X)\), and Lemma~\ref{lem:mixed-empty-positive} handles the
oriented equivalence between that anchor and the empty tuple.  In one
direction, discard the genuinely silent subderivations; in the other, reinsert
a derivation of the unique all-empty tuple at each deleted occurrence.
Induction on derivation height gives equality of all typed tuple languages.
Linearity and nondeletion are preserved for every retained child variable.
Because compression only deletes variables and never changes the relative
order of retained variables, the nonpermuting property is preserved as well.
\end{proof}

\begin{definition}[Presentation-relative characteristic sample]
\label{def:cs}
For every surviving typed nonterminal \(X\), include the identity-sector anchor
exposure
\[
  \chi(X)[\omega(X)].
\]
For every positive-rank compressed rule
\(R^+:X\to\rho(Y_1,\ldots,Y_r)\), include the identity-sector rule exposure
\[
  \chi(X)[\rho(\omega(Y_1),\ldots,\omega(Y_r))].
\]
For every compressed rank-zero rule \(R^+:X\to\tuple c\), include
\(\chi(X)[\tuple c]\).  The finite set of selected words is denoted
\(\CS(\widetilde G_0)\).
\end{definition}

\begin{lemma}[The characteristic sample is finite and positive]
\label{lem:cs-positive}
If \(L(G)\ne\emptyset\), then \(\CS(\widetilde G_0)\) is a finite nonempty
subset of \(L(G)\).
\end{lemma}

\begin{proof}
There are finitely many surviving typed symbols and compressed rules, and one
proof-side identity exposure is selected for each relevant left-hand side.
A successful start child contributes a start-child anchor exposure, so the selected set is nonempty.  Every
selected word belongs to \(L(G)\) because \(\chi(X)\) exposes all of \(L_X\)
and silent-child compression preserves the typed tuple languages.
\end{proof}

\subsection{Sample-Bounded Rank and the Canonical Learner}

Following Definition~\ref{def:tuple-occurrence}, enumerate all tuple
occurrences of arity at most \(f\) in the sample \(K\).  An
\emph{observed oriented tuple} is a pair \((\tuple x,\sigma)\) realized by
such an occurrence.  The learner uses one nonterminal
\([\tuple x;\sigma]\) for each observed oriented tuple, together with a fresh
start symbol \(\widehat S\).

\begin{definition}[Positive-support composition witness]
\label{def:composition-witness}
Let \(r\ge1\).  A \emph{positive-support composition witness of rank \(r\)}
in \(K\) consists of a parent occurrence \((E_P,\tuple z)\) of arity
\(e\le f\), with parent orientation
\(\varpi:=\operatorname{ori}(E_P)\), and an ordered \(r\)-tuple of child tuple
values \(\tuple x_j=(x_{j,1},\ldots,x_{j,d_j})\), each obtained from a sample
occurrence and satisfying \(\|\tuple x_j\|_1>0\).

Each parent component is partitioned, in order, into terminal gaps and
intervals labelled by the normalized variables
\(X_j=\{x_j^1,\ldots,x_j^{d_j}\}\).  Across all parent components every
variable is used exactly once, and the interval labelled \(x_j^k\) spells
exactly \(x_{j,k}\); coincident empty intervals retain an explicit local order.
The segmentation determines a unique linear nondeleting template
\(\rho_{\mathfrak b}\) with
\[
  \tuple z=\rho_{\mathfrak b}(\tuple x_1,\ldots,\tuple x_r).
\]
For each child put
\[
  \sigma_j:=\operatorname{ind}_j(\rho_{\mathfrak b},\varpi).
\]
The segmentation inside the parent occurrence itself realizes
\((\tuple x_j,\sigma_j)\) as an oriented tuple occurrence in the same sample
word.  Thus every state \([\tuple x_j;\sigma_j]\) used by the induced rule is
observed.
\end{definition}

\begin{example}
Let the parent tuple be \(\tuple z=(ab,cd)\), the children
\(\tuple x_1=(a,c)\) and \(\tuple x_2=(b,d)\), and suppose the parent is
exposed in orientation \(\mathrm{id}_2\).  The template
\(\rho=(x_1^1x_2^1,x_1^2x_2^2)\) induces orientation \(\mathrm{id}_2\) on both
children.  If the parent orientation is \((2,1)\), the same template induces
\((2,1)\) on both children.
\end{example}

\begin{lemma}[Sample-bounded witness rank]
\label{lem:sample-bounded-rank}
If the parent tuple \(\tuple z\) of a positive-support rank-\(r\) witness
occurs in a sample word \(w\), then
\(r\le\|\tuple z\|_1\le|w|\le\|K\|_+\).
\end{lemma}

\begin{proof}
Every child has positive total length and every child component occurs exactly
once in the nondeleting template.  Hence
\(\|\tuple z\|_1\ge\sum_j\|\tuple x_j\|_1\ge r\).  The parent components are
disjoint intervals of \(w\) apart from zero-length ties, giving the remaining
inequalities.
\end{proof}

\begin{lemma}[Enumeration of arbitrary-rank witnesses]
\label{lem:general-witness-enumeration}
Fix \(f\) and a finite sample \(K\), and put \(n:=\|K\|_+\).  All
positive-support composition witnesses, together with parent and induced child
orientations and canonical rule encodings, can be effectively enumerated.  For
every fixed branching cap \(b\), all witnesses of rank at most \(b\) are
enumerable in \(n^{O(fb)}\) time and space.
\end{lemma}

\begin{proof}
An arity-\(d\) occurrence is determined by \(2d\) cuts and one of at most
\(d!\le f!\) orientations, so the enumeration of oriented occurrences differs
only by a factor depending on the fixed fan-out bound \(f\).  For a rank-\(r\)
candidate,
enumerate parent and child values, variable intervals, terminal gaps, and local
orders of coincident empty intervals.  There are \(n^{O(fr)}\) candidates.
The induced child orientations are then read deterministically from the
resulting template and parent orientation.  Lemma~\ref{lem:sample-bounded-rank}
bounds the uncapped search by \(r\le n\).
\end{proof}

In particular, for fixed \(f\), all positive-support witnesses of rank at
most two are enumerable in \(\|K\|_+^{O(f)}\) time and space.

\begin{lemma}[Identity-exposed compressed rules occur in the enumeration]
\label{lem:exposed-rule-enumerated}
Let \(R^+:X\to\rho(Y_1,\ldots,Y_r)\) be a positive-rank rule of
\(\widetilde G_0^+\), and put
\[
  \tuple z_R:=\rho(\omega(Y_1),\ldots,\omega(Y_r)).
\]
If the characteristic sample is contained in \(K\), the witness enumeration
contains the rule
\[
  [\tuple z_R;\mathrm{id}_{\mu(X)}]\to
  \rho([\omega(Y_1);\mathrm{id}_{\mu(Y_1)}],\ldots,
       [\omega(Y_r);\mathrm{id}_{\mu(Y_r)}]).
\]
\end{lemma}

\begin{proof}
The rule exposure \(\chi(X)[\tuple z_R]\) supplies the parent occurrence in
identity orientation.  Its ground segmentation by the variables of \(\rho\)
realizes each child anchor inside the same sample word.  By
Lemma~\ref{lem:identity-orientation-propagation}, every retained child is also
realized in identity orientation.  Every retained child is non-silent and
hence has positive total length.  The segmentation therefore satisfies
Definition~\ref{def:composition-witness}, and completeness of the witness
enumeration emits the displayed rule.
\end{proof}

\begin{definition}[Capped and general orientation-aware canonical hypotheses]
\label{def:learner}
For \(b\ge1\), the capped hypothesis \(\widehat G_{\le b}(K)\) has start
symbol \(\widehat S\), all nonterminals \([\tuple x;\sigma]\) realized by
oriented tuple occurrences in \(K\), and the following rules:
\begin{enumerate}[label=(\roman*),leftmargin=*]
\item for every \(w\in K\),
      \(\widehat S\to[(w);\mathrm{id}_1]\);
\item for every observed oriented tuple, the constant rule
      \([\tuple x;\sigma]\to\tuple x\);
\item for every positive-support witness of rank \(1\le r\le b\), with
      parent orientation \(\varpi\), parent tuple \(\tuple z\), template \(\rho\),
      and induced child orientations
      \(\sigma_j=\operatorname{ind}_j(\rho,\varpi)\), the rule
      \[
        [\tuple z;\varpi]\to
        \rho([\tuple x_1;\sigma_1],\ldots,[\tuple x_r;\sigma_r]);
      \]
\item for observed \((\tuple x,\sigma)\) and \((\tuple y,\sigma)\) of the
      same arity, the identity-template rule
      \([\tuple x;\sigma]\to[\tuple y;\sigma]\) whenever
      \(h^{(d)}(\tuple x)=h^{(d)}(\tuple y)\) and some concrete context \(E\)
      of orientation \(\sigma\) satisfies
      \(E[\tuple x],E[\tuple y]\in K\).
\end{enumerate}
If \(K\ne\emptyset\), define
\(\widehat G(K):=\widehat G_{\le\|K\|_+}(K)\); for \(K=\emptyset\), use the
grammar with no productive start rule.  Fixed encoding orders make the learner
deterministic and set-driven.  Its only inputs are \(K,f,h\); orientations are
read from sample occurrences.
\end{definition}

\begin{remark}[Compressing semantic unary rules]
The definition may add semantic unary rules for many observed pairs.  An
implementation may first partition observed oriented tuples of each arity, orientation, and
\(h\)-type by the shared-context relation generated inside that sector and choose
one representative per resulting class, connecting other members only to that
representative.  This reduces redundant unary rules without changing the
language generated by the hypothesis or any complexity claim below.
\end{remark}

\begin{lemma}[Finiteness and effectiveness of the canonical hypotheses]
\label{lem:learner-effective}
For every finite \(K\) and every \(b\ge1\), both
\(\widehat G_{\le b}(K)\) and \(\widehat G(K)\) are finite standard MCFGs of
fan-out at most \(f\).  They are effectively constructible from \(K,f,h\);
when \(K\ne\emptyset\), every rule of \(\widehat G(K)\) has rank at most
\(\|K\|_+\).
\end{lemma}

\begin{proof}
A finite sample has only finitely many tuple occurrences of every arity at most
\(f\), hence only finitely many observed tuple values and semantic unary-rule
candidates.  In rule~(iv), every relevant concrete context is obtained by
replacing a marked tuple occurrence in one of the finitely many sample words by
named holes, so only finitely many candidate contexts need be inspected.
Lemma~\ref{lem:general-witness-enumeration} effectively enumerates
all witness rules below any fixed cap, and
Lemma~\ref{lem:sample-bounded-rank} makes the sample-dependent cap
\(\|K\|_+\) sufficient for the general hypothesis.  The start and constant
rules are finite and effective by inspection.  Every nonterminal
\([\tuple x;\sigma]\) has fan-out equal to the arity of \(\tuple x\), at most
\(f\).  The orientation is a finite state annotation and does not alter the
derived tuple.  Every emitted composition template is linear and nondeleting,
hence a standard MCFG rule.
\end{proof}


\paragraph{Size of the selected witness set.}
If a proof-side presentation has nonterminal set $V$, rule set $P$, fan-out at
most $f$, and rule rank at most $b$, the output-type refinement has at most
$|V||M|^f$ typed nonterminals and $|P||M|^{fb}$ typed rule instances.  Thus
the characteristic sample selects at most that many state and rule exposures.
For a fixed finite proof-side grammar, these witnesses are effectively
computable: enumerate successful derivation trees in dovetailing fashion and
retain the first required state or rule exposure when it appears.  This counts
selected words only and does not give a polynomial bound on their total
lengths.

\subsection{Soundness and Exact Reconstruction}
\subsubsection{Soundness}
\label{sec:soundness}

\begin{lemma}[Induced child context and its orientation]
\label{lem:filling-identity}
Let \(\rho\) be a rank-\(r\) linear nondeleting template, let \(E\) be a
sentence context for its output tuple with orientation \(\varpi\), and fix
compatible tuples for all children except child \(j\).  Replacing every
variable of child \(j\) by its named hole and every other child variable by
the corresponding fixed component produces a sentence context \(E_j\) such
that
\[
  E_j[\tuple u_j]=E[\rho(\tuple u_1,\ldots,\tuple u_r)]
\]
and
\[
  \operatorname{ori}(E_j)=\operatorname{ind}_j(\rho,\varpi).
\]
The identity and orientation remain valid after arbitrary replacements of the
fixed sibling tuples.
\end{lemma}

\begin{proof}
Linearity and nondeletion make every variable of child \(j\) occur exactly
once.  Scanning the parent components in the order \(\varpi\) gives precisely the
variable order used in Definition~\ref{def:induced-orientation}.  Filling the
holes is composition of the child substitution with the template and then
with \(E\).
\end{proof}

\begin{lemma}[Witnessed oriented composition preserves equivalence]
\label{lem:witnessed-composition}
Let \(L\) be \((f,h)\)-tuple-substitutable, let \(\rho\) be a rank-\(r\)
linear nondeleting template, and let \(E\) be an accepting parent context of
orientation \(\varpi\):
\[
  E[\rho(\tuple x_1,\ldots,\tuple x_r)]\in L.
\]
Put \(\sigma_j:=\operatorname{ind}_j(\rho,\varpi)\).  If
\[
  \tuple u_j\equiv_{L,h,\sigma_j}\tuple x_j
  \qquad(1\le j\le r),
\]
then
\[
  \rho(\tuple u_1,\ldots,\tuple u_r)
  \equiv_{L,h,\varpi}
  \rho(\tuple x_1,\ldots,\tuple x_r).
\]
\end{lemma}

\begin{proof}
Replace children one at a time.  Lemma~\ref{lem:filling-identity} gives, at
stage \(j\), an induced accepting context of orientation \(\sigma_j\).
Equality of the corresponding sector distributions permits replacement of
\(\tuple x_j\) by \(\tuple u_j\).  After all replacements the parent remains
accepted by \(E\).  Homomorphic template evaluation and equality of child
\(h\)-types give equality of the parent \(h\)-types.  The original and updated
parents share the accepting context \(E\) in sector \(\varpi\), so
Lemma~\ref{lem:shared-context} gives the displayed parent equivalence.
\end{proof}

\begin{proposition}[Soundness of every capped hypothesis]
\label{prop:learner-soundness}
Let \(L\) be \((f,h)\)-tuple-substitutable and \(K\subseteq L\).  For every
\(b\ge1\), if
\[
  [\tuple x;\sigma]\Rightarrow^*\tuple u
\]
in \(\widehat G_{\le b}(K)\), then
\(\tuple u\equiv_{L,h,\sigma}\tuple x\).  Consequently
\[
  L(\widehat G_{\le b}(K))\subseteq L,
  \qquad
  L(\widehat G(K))\subseteq L.
\]
\end{proposition}

\begin{proof}
Induct on the derivation.  A constant rule is reflexive.  A semantic unary
rule is sound by Lemma~\ref{lem:shared-context}, and all such rules preserve
the same orientation.  For a witness rule, its parent occurrence supplies an
accepted parent application in sector \(\varpi\); apply the induction hypotheses
to the children and Lemma~\ref{lem:witnessed-composition}.  Finally every
start derivation begins with \(\widehat S\to[(w_0);\mathrm{id}_1]\) for some
\(w_0\in K\).  The empty unary context then places the derived word in \(L\).
\end{proof}

\subsubsection{Completeness and Exact Reconstruction}
\label{sec:completeness}

\begin{proposition}[Completeness for a fixed normalized presentation]
\label{prop:completeness}
Let \(G\) be a finite nondeleting, nonpermuting MCFG of fan-out at most \(f\),
let \(L=L(G)\) be \((f,h)\)-tuple-substitutable, and construct
\(\widetilde G_0^+\) and \(\CS(\widetilde G_0)\) as above.  Let \(b_G\) be
the maximum rank of a compressed rule.  If
\[
  \CS(\widetilde G_0)\subseteq K\subseteq L,
\]
then for every \(b\ge\max\{1,b_G\}\),
\[
  L\subseteq L(\widehat G_{\le b}(K)),
\]
and also \(L\subseteq L(\widehat G(K))\).
\end{proposition}

\begin{proof}
For every surviving typed nonterminal \(X\), put
\[
  \iota(X):=[\omega(X);\mathrm{id}_{\mu(X)}].
\]
If \(\widetilde S\to X\) survives, then \(\chi(X)=\blank_1\), so the learner's
start rule to \(\iota(X)\) is present.

Fix a compressed rank-zero rule \(R:X\to\tuple c\).  The anchor and rule
exposures share the identity context \(\chi(X)\), and \(\tuple c\) has the
output type of \(X\).  Hence the semantic unary rule
\[
  \iota(X)\to[\tuple c;\mathrm{id}_{\mu(X)}]
\]
is present, followed by the constant rule.  If \(X\) can derive both the
all-empty tuple and positive tuples, the required identity-sector replacement
is justified by Lemma~\ref{lem:mixed-empty-positive}.

Now fix a positive-rank compressed rule
\(R:X\to\rho(Y_1,\ldots,Y_r)\), and put
\[
  \tuple z_R:=\rho(\omega(Y_1),\ldots,\omega(Y_r)).
\]
The parent anchor and rule exposure share \(\chi(X)\) and have the same output
type, so the semantic unary rule
\[
  \iota(X)\to[\tuple z_R;\mathrm{id}_{\mu(X)}]
\]
is present.  By Lemma~\ref{lem:exposed-rule-enumerated}, the learner also
contains
\[
  [\tuple z_R;\mathrm{id}_{\mu(X)}]\to
  \rho(\iota(Y_1),\ldots,\iota(Y_r))
\]
whenever its cap is at least \(r\).  Induction on a successful
\(\widetilde G_0^+\)-derivation therefore simulates the whole normalized
derivation using only identity-oriented proof-side states.  This restriction
is purely a completeness device: Proposition~\ref{prop:learner-soundness}
still requires the learner to maintain arbitrary orientation sectors for
sample-derived rules.  At the root the unique orientation is
\(\mathrm{id}_1\), giving every word of \(L\).

For the uncapped learner, every required witness has a rule exposure word
\(w_R\in K\), and Lemma~\ref{lem:sample-bounded-rank} gives
\(r\le|w_R|\le\|K\|_+\).  Hence all required witnesses lie below the
sample-dependent cap.
\end{proof}

\begin{theorem}[Exact reconstruction for fixed-observation MCFGs]
\label{thm:exact-reconstruction}
Fix \(f\ge1\) and an explicit finite homomorphism \(h:\Sigma^*\to M\).  Let
\[
  L\in\Cmcf{f}{h}.
\]
Then there exists a finite characteristic sample \(S_L\subseteq L\) such that,
for every finite
\[
  S_L\subseteq K\subseteq L,
\]
the general canonical learner satisfies
\[
  L(\widehat G(K))=L.
\]
If \(L\ne\emptyset\), the sample can be chosen as
\(S_L=\CS(\widetilde G_0)\) for a nondeleting, nonpermuting proof-side
presentation of \(L\).  For \(L=\emptyset\), take \(S_L=\emptyset\).
\end{theorem}

\begin{proof}
If \(L=\emptyset\), then \(K=\emptyset\) is the only finite set satisfying
\(K\subseteq L\), and \(\widehat G(\emptyset)\) has empty language by
definition.  Assume \(L\ne\emptyset\).  By Definition~\ref{def:class-cmcf}, choose a finite
standard \(f\)-MCFG \(G\) with \(L(G)=L\).  Since \(G\) is finite it has a
finite branching bound.  By the rank-preserving normalization recorded by
Kanazawa \cite[Section~2]{Kanazawa2019Ogden}, choose an equivalent nondeleting,
nonpermuting \(f\)-MCFG \(G_{\mathrm{np}}\) with the same branching bound.  In the
explicit normalization used above, this is realized by deletion elimination
\cite[Lemma~2.2]{SekiEtAl1991} followed by
Lemma~\ref{lem:permutation-decoration}.  Apply the refinement, compression, and characteristic-sample
construction above to \(G_{\mathrm{np}}\).  Proposition~\ref{prop:completeness}
gives
\[
  L\subseteq L(\widehat G(K)),
\]
while Proposition~\ref{prop:learner-soundness} gives the reverse inclusion.
\end{proof}

\begin{corollary}[TxtBc identification from positive data]
\label{cor:semantic-identification}
For every fixed \(f\) and fixed explicit \(h\), the full semantic slice
\(\Cmcf{f}{h}\) is TxtBc-identifiable from positive data by the set-driven
learner
\[
  K\longmapsto\widehat G(K).
\]
No target-specific rule-rank bound or reduced-presentation promise is part of the target class.
\end{corollary}

\begin{proof}
For every nonempty target, the content of every text eventually contains the
finite characteristic sample from Theorem~\ref{thm:exact-reconstruction},
after which all hypothesis languages equal the target.  For the empty target,
every text consists entirely of pauses, so its content is always empty and the
learner always returns the empty language.
\end{proof}

\begin{corollary}[TxtEx identification by a conservative wrapper]
\label{cor:gold-explanatory}
For every fixed \(f\) and fixed explicit \(h\), the full semantic slice
\(\Cmcf{f}{h}\) is TxtEx-identifiable from positive data: there is a
computable learner whose hypothesis grammar itself is eventually constant and
generates the target language.
\end{corollary}

\begin{proof}
Run a conservative wrapper around the canonical learner.  Initialize the
current hypothesis as \(H:=\widehat G(\emptyset)\).  Pauses are ignored.  When a positive example \(w\)
is received, keep \(H\) unchanged if \(w\in L(H)\); otherwise recompute
\(H:=\widehat G(K)\) from the current finite content \(K\).  Membership in a
fixed finite MCFG is decidable, so this wrapper is computable.

Throughout the run, the current hypothesis contains every positive datum seen so
far.  Indeed, after a recomputation this follows from
\(K\subseteq L(\widehat G(K))\), and when \(H\) is kept the newly seen word
already lies in \(L(H)\).

For a target \(L\in\Cmcf{f}{h}\), every recomputed canonical hypothesis is
sound by Proposition~\ref{prop:learner-soundness}, hence its language is a
subset of \(L\).  If the current hypothesis language is a proper subset of
\(L\), some missing target word eventually appears on every text and forces a
recomputation.  Once the finite characteristic sample of
Theorem~\ref{thm:exact-reconstruction} is contained in the observed content,
the next recomputation, if any, produces a grammar whose language is exactly
\(L\).  If no recomputation occurs after that point, the current hypothesis is
still sound and accepts every later positive datum; since every word of $L$
eventually appears on the text, it already contains all of $L$ and hence is
also exact.  From the first exact stage onward every positive example is
accepted and the wrapper never changes that grammar again.  If a correct
grammar is reached earlier, stabilization occurs earlier.  The empty target is
immediate.
\end{proof}


\section{Effective Hypothesis Construction}
\label{sec:poly-time}

The qualitative learner allows arbitrary finite rule rank and is therefore not
claimed to run in polynomial time on a finite sample; throughout this section
we use the slicewise convention of Remark~\ref{rem:slicewise-poly}.  This section separates
that issue from TxtBc-identifiability.  If the composition rank is capped by a fixed \(b\), the same witness
construction is polynomial.  Characteristic-data size is a distinct question
and can remain exponential even when hypothesis construction is polynomial.

\paragraph{Encoding convention.}
We use canonical finite encodings for all sample-derived objects.  An occurrence is
recorded by its sample-word index, component permutation, and cut positions;
tuple contexts and templates are recorded by arity together with delimited
terminal factors and ordered hole or variable labels.

For an explicitly supplied observation \(h:\Sigma^*\to M\), let
\(\|h\|_{\mathrm{enc}}\) denote the bit length of its multiplication table and
the letter images.  The first theorem below is slicewise in fixed \(f,b,h\).

\begin{theorem}[Slicewise-polynomial construction under a branching cap]
\label{thm:poly}
Fix a fan-out bound \(f\ge1\), a branching cap \(b\ge1\), and an explicit
finite-monoid homomorphism \(h:\Sigma^*\to M\).  From any finite positive
sample \(K\), the capped canonical hypothesis
\[
  \widehat G_{\le b}(K)
\]
can be constructed, including its complete output encoding, in time
\[
  (1+\|K\|_+)^{O(fb)}.
\]
Its fan-out is at most \(f\), and every composition-witness rule has rank at
most \(b\) (semantic unary rules and rank-zero constant rules are also
allowed).  The polynomial degree may depend on \(f\) and \(b\).
\end{theorem}

\begin{proof}
Put \(n:=\max(1,\|K\|_+)\).  Tuple occurrences of arity at most \(f\) are
enumerable in \(n^{O(f)}\) time.%
\footnote{An arity-\(d\) occurrence in a word of length \(m\) is specified by
two cut positions for each component together with the order of the named
components, so the number of candidates is at most
\((m+1)^{2d}d!\).  Moreover \(|K|\le\|K\|_+\le n\) and
\(|w|\le n\) for every \(w\in K\).  Hence all occurrences with
\(d\le f\) can be enumerated in \(n^{O(f)}\) time for fixed \(f\).}  All observed constant rules are read from
these tuples.  Semantic unary rules are obtained by grouping occurrences by
arity and concrete-context encoding and comparing componentwise \(h\)-types;
even exhaustive pairwise comparison remains within \(n^{O(f)}\).

By Lemma~\ref{lem:general-witness-enumeration}, all positive-support composition
witnesses of rank at most \(b\), together with their canonical template
encodings, are enumerable in \(n^{O(fb)}\) time and space.  Canonical sorting
removes duplicate nonterminals and rules within the same bound.  Thus the
number of output symbols, rules, and their total encoding length are all
\(n^{O(fb)}\).
\end{proof}


\paragraph{Scope of the bound.}
The theorem concerns construction of the current capped hypothesis, including
its output encoding, from a supplied finite sample.  It does not give a
polynomial bound on the total length of a characteristic sample.  The general
learner uses the sample-dependent cap $\|K\|_+$ and is therefore a qualitative
identification result rather than a polynomial-time algorithm with an unknown
rank parameter.


\section{Learning with Latent Observations}
\label{sec:fixed-observation}

The reconstruction theorem assumes that the finite observation is supplied to
the learner.  We now weaken that advice.  First we assume only an upper bound
on the observation size and compile all compatible finite observations into one
universal refinement.  We then quantify the positive information that any
learner may be forced to see on such bounded slices, and finally remove the
size bound altogether.

\subsection{Bounded Latent Observations}
\begin{definition}[Bounded-size observation union]
\label{def:bounded-observation-union}
Fix the finite alphabet \(\Sigma\), a fan-out bound \(f\ge1\), and \(k\ge1\).
Define
\[
  \Clatent{f}{k}:=\bigcup_{\substack{h:\Sigma^*\to M\\ |M|\le k}} \Cmcf{f}{h},
\]
where the union ranges over all finite monoid homomorphisms
\(h:\Sigma^*\to M\) with \(|M|\le k\).
\end{definition}

\begin{remark}[Embedding invariance]
If \(\varphi:M\to M'\) is injective, then
\(\Cmcf{f}{h}=\Cmcf{f}{\varphi\circ h}\), since
Definition~\ref{def:tuple-substitutable} uses the observation only through
equalities of componentwise values.
\end{remark}

\begin{definition}[Universal bounded observation]
\label{def:universal-bounded-observation}
Fix a nonempty finite alphabet \(\Sigma\) and \(k\ge1\).  An explicit finite
homomorphism
\[
  U:\Sigma^*\to N
\]
is \emph{\(k\)-universal} if every explicit homomorphism
\(h:\Sigma^*\to M\) with \(|M|\le k\) factors through \(U\); equivalently,
\(h\preceq U\) for every such \(h\).
\end{definition}

\begin{proposition}[Bounded observation size restores TxtBc-identifiability]
\label{prop:bounded-union-identifiable}
For every fixed finite alphabet \(\Sigma\), fan-out bound \(f\ge1\), and
\(k\ge1\), there is an explicit finite monoid homomorphism
\[
  U_{\Sigma,k}:\Sigma^*\twoheadrightarrow M_{\Sigma,k}^{\mathrm{univ}}
\]
computable from \((\Sigma,k)\), such that
\[
  \Clatent{f}{k}\subseteq\Cmcf{f}{U_{\Sigma,k}}.
\]
Every explicit homomorphism \(h:\Sigma^*\to M\) with \(|M|\le k\) factors
through \(U_{\Sigma,k}\).  Consequently the canonical learner with parameters
\((f,U_{\Sigma,k})\) TxtBc-identifies \(\Clatent{f}{k}\) from positive data.
\end{proposition}

\begin{proof}
The construction can be made completely effective without choosing an
abstract system of isomorphism representatives.  For each
\(1\le j\le k\), enumerate all binary operation tables on the standard carrier
\([j]=\{1,\ldots,j\}\), retaining exactly those satisfying the monoid axioms.
For every retained monoid table, enumerate all maps \(\Sigma\to[j]\); by
freeness, each such map determines a unique homomorphism from \(\Sigma^*\).
This produces a finite effective list
\(h_i:\Sigma^*\to M_i\), \(1\le i\le r\), containing every homomorphism into
a monoid of size at most \(k\) up to relabeling of its carrier.

Form the temporary product
\[
  P_{\Sigma,k}:=M_1\times\cdots\times M_r,
  \qquad
  H^{\mathrm{prod}}_{\Sigma,k}:=(h_1,\ldots,h_r):
     \Sigma^*\to P_{\Sigma,k}.
\]
The product is used only to make finiteness and effectivity transparent.
Compute the submonoid generated by the finitely many letter values
\(H^{\mathrm{prod}}_{\Sigma,k}(a)\), \(a\in\Sigma\), and put
\[
  M_{\Sigma,k}^{\mathrm{univ}}
  :=\operatorname{im}(H^{\mathrm{prod}}_{\Sigma,k}).
\]
Let
\[
  U_{\Sigma,k}:\Sigma^*\twoheadrightarrow M_{\Sigma,k}^{\mathrm{univ}}
\]
be the same product map with codomain restricted to its image.  Its
multiplication table and its letter values are computable from \((\Sigma,k)\),
so \(U_{\Sigma,k}\) is an explicit finite observation.

Let \(h:\Sigma^*\to M\) be explicit with \(|M|=j\le k\).  Choose a
bijection \(\varphi:M\to[j]\) and transport the multiplication table of \(M\)
along \(\varphi\).  The resulting table and the letter map
\(\varphi\circ h|_\Sigma\) occur in the enumeration, so
\(\varphi\circ h=h_i\) for some coordinate \(i\).  Let \(\operatorname{pr}_i\) be that coordinate projection.  Then
\(\pi_h:=\varphi^{-1}\circ\operatorname{pr}_i|_{M_{\Sigma,k}^{\mathrm{univ}}}\)
is a monoid homomorphism \(M_{\Sigma,k}^{\mathrm{univ}}\to M\) with
\[
  h=\pi_h\circ U_{\Sigma,k}.
\]
Hence \(h\preceq U_{\Sigma,k}\), and
Proposition~\ref{prop:h-refinement-monotonicity} gives
\(\Cmcf{f}{h}\subseteq\Cmcf{f}{U_{\Sigma,k}}\).  Taking the union over all
\(|M|\le k\) yields the displayed inclusion.  By
Corollary~\ref{cor:semantic-identification}, the canonical learner for the
fixed morphism \(U_{\Sigma,k}\) TxtBc-identifies
\(\Cmcf{f}{U_{\Sigma,k}}\), and hence its bounded-observation subclass.
\end{proof}


\subsection{The Information Cost of Universalization}
\label{subsec:uniform-observation-bound}

For learner-uniform lower bounds we use the standard semantic locking notion.
\begin{definition}[Semantic TxtBc-locking sequence]
\label{def:semantic-locking-sequence}
Let \(\mathcal A\) be an arbitrary learner and let \(L\subseteq\Sigma^*\).
A finite sequence
\[
  \sigma\in(L\cup\{\#\})^*
\]
is a \emph{semantic TxtBc-locking sequence for \(\mathcal A\) on \(L\)} if,
for every finite continuation \(\tau\in(L\cup\{\#\})^*\),
\[
  L(\mathcal A(\sigma\tau))=L.
\]
Its \emph{locking content} is \(\operatorname{content}(\sigma)\), and its
positive locking-content size is
\(\|\operatorname{content}(\sigma)\|_+\).
Thus the certificate ignores order and multiplicity only when measuring its
positive information content; the learner itself need not be set-driven.
\end{definition}

\begin{lemma}[Semantic locking lemma]
\label{lem:semantic-locking-exists}
If an arbitrary learner \(\mathcal A\) TxtBc-identifies a nonempty language
\(L\) from positive data, then \(\mathcal A\) has a semantic TxtBc-locking
sequence on \(L\).
\end{lemma}

\begin{proof}
Suppose no such sequence exists.  Choose a sequence
\(x_0,x_1,x_2,\ldots\) over \(L\) in which every member of \(L\) occurs at
least once; if \(L\) is finite, repeat its elements thereafter.  Starting with
\(\sigma_0=\eps\), construct finite prefixes inductively.  Given
\(\sigma_i\), put \(\rho_i=\sigma_i x_i\).  Since \(\rho_i\) is not a
locking sequence, there is a finite
\(\tau_i\in(L\cup\{\#\})^*\) such that
\[
  L(\mathcal A(\rho_i\tau_i))\ne L.
\]
Set \(\sigma_{i+1}=\rho_i\tau_i\).  The infinite concatenation
\[
  x_0\tau_0x_1\tau_1x_2\tau_2\cdots
\]
is a text for \(L\), but at the end of every block \(x_i\tau_i\) the
learner is semantically incorrect.  Hence the learner is incorrect infinitely
often on this text, contradicting TxtBc-identification.
\end{proof}


\begin{lemma}[Fixed-length regular slices]
\label{lem:fixed-length-regular-slices}
Let \(N\ge0\), let \(g:\Sigma^*\to M\) be a finite monoid homomorphism, and
let \(P\subseteq M\).  Then
\[
  L:=\Sigma^N\cap g^{-1}(P)
\]
is \((f,g)\)-tuple-substitutable for every \(f\ge1\).
\end{lemma}

\begin{proof}
Suppose \(\tuple x,\tuple y\) have equal componentwise \(g\)-type and share
an accepting sentence context \(E\).  Since both fillings have length \(N\),
the total lengths of the two tuples are equal.  Hence in every other sentence
context \(F\), the two fillings again have the same total length.  Their
\(g\)-values are also equal, because corresponding components have equal
\(g\)-values.  Thus either both fillings or neither belongs to
\(\Sigma^N\cap g^{-1}(P)\), in every orientation sector.
\end{proof}


\begin{lemma}[Fixed-length deletion family]
\label{lem:fixed-length-deletion-family}
Let \(\Gamma=\{0,1\}\), let \(n\ge1\), and put
\[
  F_n:=\Gamma^n,\qquad
  F_{n,w}:=\Gamma^n\setminus\{w\}\quad(w\in\Gamma^n),
  \qquad
  k_n:=\frac{n(n+1)}2+2.
\]
For every fan-out bound \(f\ge1\),
\[
  F_n\in\Clatent{f}{1},
  \qquad
  F_{n,w}\in\Clatent{f}{k_n}
  \quad(w\in F_n).
\]
\end{lemma}

\begin{proof}
The language \(F_n\) is a fixed-length regular slice and is witnessed by the
one-element observation, so $F_n\in\Clatent{f}{1}$.  For a fixed
\(w\in\Gamma^n\), let
\(\eta_w:\Gamma^*\to M_w\) be the syntactic morphism of the singleton
language \(\{w\}\).  There are at most \(n(n+1)/2\) nonempty factors of
\(w\), one additional class containing \(\eps\), and one junk class containing
all nonfactors.  Hence
\[
  |M_w|\le \frac{n(n+1)}2+2=k_n.
\]
Moreover
\[
  F_{n,w}=\Gamma^n\cap
  \eta_w^{-1}\bigl(M_w\setminus\{\eta_w(w)\}\bigr),
\]
because the empty two-sided context already distinguishes \(w\) from every
other length-\(n\) word.  The fixed-length regular-slice lemma therefore gives
\((f,\eta_w)\)-tuple substitutability, proving
\(F_{n,w}\in\Clatent{f}{k_n}\).
\end{proof}

\begin{lemma}[Deletion-family locking obstruction]
\label{lem:deletion-locking-obstruction}
Let \(L\ne\emptyset\) and let \(\mathcal A\) TxtBc-identify \(L\) and every
one-point deletion \(L\setminus\{w\}\), \(w\in L\).  Then every semantic
TxtBc-locking sequence \(\sigma\) for \(\mathcal A\) on \(L\) satisfies
\[
  \operatorname{content}(\sigma)=L.
\]
\end{lemma}

\begin{proof}
If some \(w\in L\) were missing from the locking content, the locking prefix
could be extended to a text for \(L\setminus\{w\}\).  Locking would force all
subsequent hypotheses to denote \(L\), whereas TxtBc-identification of the
deletion language would eventually force them to denote
\(L\setminus\{w\}\), a contradiction.
\end{proof}

\begin{lemma}[Locking content and characteristic samples]
\label{lem:locking-characteristic-bridge}
Let \(\mathcal A\) be set-driven and let \(L\subseteq\Sigma^*\).
A finite set \(S\subseteq L\) is characteristic for \(\mathcal A\) on \(L\)
if and only if any finite sequence over \(L\cup\{\#\}\) with content \(S\)
is a semantic TxtBc-locking sequence.
\end{lemma}

\begin{proof}
If \(S\) is characteristic, every continuation has finite content
\(K\) with \(S\subseteq K\subseteq L\), so set-drivenness and the
characteristic property keep the hypothesis language equal to \(L\).
Conversely, if a sequence with content \(S\) is locking and
\(S\subseteq K\subseteq L\) is finite, append exactly the elements of
\(K\setminus S\); set-drivenness then gives the characteristic property.
\end{proof}

The deletion-family locking step below is the standard positive-data trap:
a learner locked on a finite language cannot omit a word while also learning
the corresponding one-point deletion.  The quantitative content specific to
the present setting is Lemma~\ref{lem:fixed-length-deletion-family}, which
places every deletion inside an observation slice of size only $O(n^2)$.

\begin{theorem}[Learner-uniform \(2^{\Omega(\sqrt{k})}\) locking lower bound]
\label{thm:no-uniform-polynomial-k-data}
Fix \(\Gamma=\{0,1\}\) and \(f\ge1\), and put
\(k_n=n(n+1)/2+2\).  For every \(n\) and every arbitrary learner
\(\mathcal A_{k_n}\) that TxtBc-identifies the full semantic slice
\(\Clatent{f}{k_n}\), every semantic locking sequence \(\sigma\) for the
target \(F_n=\Gamma^n\) satisfies
\[
  \operatorname{content}(\sigma)=F_n,
  \qquad
  \|\operatorname{content}(\sigma)\|_+=n2^n.
\]
Thus the necessary positive locking information is
\(2^{\Omega(\sqrt{k_n})}\).  If the learner is set-driven, every
characteristic sample for \(F_n\) is exactly \(F_n\).
\end{theorem}

\begin{proof}
Existence of a semantic locking sequence for each correctly TxtBc-learned target
is guaranteed by Lemma~\ref{lem:semantic-locking-exists}.  Lemma~\ref{lem:fixed-length-deletion-family} puts \(F_n\) and all one-point
deletions \(F_{n,w}\) in the same semantic slice
\(\Clatent{f}{k_n}\).  Lemma~\ref{lem:deletion-locking-obstruction} therefore
forces every locking content to equal \(F_n\), whence its positive size is
\[
  \sum_{w\in\Gamma^n}|w|=n2^n.
\]
Since \(k_n=\Theta(n^2)\), the stated asymptotic form follows.  The set-driven
statement is Lemma~\ref{lem:locking-characteristic-bridge}.
\end{proof}

\begin{corollary}[Lower bound for arbitrary observation budgets]
\label{cor:locking-lower-bound-all-k}
For every sufficiently large integer $k$, every learner that TxtBc-identifies
$\Clatent{f}{k}$ has a target in that slice whose every semantic locking
sequence has positive content of total length $2^{\Omega(\sqrt{k})}$.
\end{corollary}

\begin{proof}
Choose $n$ maximal with $k_n=n(n+1)/2+2\le k$.  Then
$\Clatent{f}{k_n}\subseteq\Clatent{f}{k}$, so
Theorem~\ref{thm:no-uniform-polynomial-k-data} applies to the target $F_n$.
Maximality gives $k<k_{n+1}=\Theta(n^2)$, hence
$n=\Theta(\sqrt{k})$, and the required locking content has total length
$n2^n=2^{\Omega(\sqrt{k})}$.
\end{proof}


\subsection{Unbounded Latent Observations}
\begin{definition}[Unbounded latent-observation union]
\label{def:no-advice-union}
Fix an alphabet \(\Sigma\) and a fan-out bound \(f\).  We now allow the target
to belong to any fixed-observation slice, but we do not tell the learner which
finite morphism witnesses this membership.  For this reason, define the
\emph{unbounded latent-observation union}
\[
  \Callobs{f}
  :=
  \bigcup_h \Cmcf{f}{h},
\]
where the union ranges over all explicit finite monoid homomorphisms
\[
  h:\Sigma^*\to M.
\]
A learner for \(\Callobs{f}\) receives only positive data and the fixed
fan-out bound; in particular, it receives neither a witnessing morphism nor a
bound on the size of its codomain monoid.
\end{definition}

Proposition~\ref{prop:copy-infinite-observation} shows that COPY belongs to no
finite-observation slice.  Hence
$\Callobs{2}\subsetneq2\text{-}\mathrm{MCFL}$.

\begin{lemma}[Ascending-chain obstruction]
\label{lem:ascending-chain-obstruction}
Let a language class \(\mathcal C\) contain languages
\[
  L_1\subsetneq L_2\subsetneq\cdots
\]
and their union \(L_\infty=\bigcup_{k\ge1}L_k\).  Then \(\mathcal C\) is not TxtBc-identifiable from positive data.
\end{lemma}

\begin{proof}
Assume that a learner \(\mathcal A\) TxtBc-identifies every language in
\(\mathcal C\).  We construct a text for \(L_\infty\) on which
\(\mathcal A\) outputs infinitely many different hypothesis languages.
Fix an enumeration \(v_1,v_2,\ldots\) of \(L_\infty\).  Suppose a finite prefix
\(\sigma_s\) has been constructed and its content is contained in some
\(L_{k_s}\).  Continue \(\sigma_s\) with elements of a text for \(L_{k_s}\).
The resulting infinite continuation is a text for \(L_{k_s}\), so at some
finite extension the learner must output a hypothesis whose language is
\(L_{k_s}\).  Then append an element of
\(L_{k_s+1}\setminus L_{k_s}\), and append \(v_s\) if it has not yet appeared.
The new finite content is contained in some later \(L_{k_{s+1}}\), and the
construction can be repeated.

Every \(v_s\) eventually appears, so the resulting sequence is a text for
\(L_\infty\).  At the selected stages, however, the learner outputs the
strictly increasing approximant languages \(L_{k_s}\), and therefore cannot
converge semantically on this text.  This contradicts TxtBc-identification of
\(L_\infty\).
\end{proof}

\begin{corollary}[Regular languages lie in the latent-observation union]
\label{cor:regular-in-latent-union}
For every \(f\ge1\),
\[
  \mathrm{REG}\subseteq\Callobs{f}.
\]
\end{corollary}

\begin{proof}
For a regular language $L$, let $\eta_L$ be its finite syntactic morphism and
write $L=\eta_L^{-1}(P)$ for the corresponding accepting set $P$.  Apply
Proposition~\ref{prop:h-recognizable-in-slice}.
\end{proof}

The following is the standard superfinite obstruction underlying Gold-style
and finite-tell-tale impossibility arguments.

\begin{theorem}[Failure of TxtBc-identifiability with unbounded latent observation]
\label{thm:no-advice-nonidentifiability}
Fix a nonempty finite alphabet \(\Sigma\).  For every fixed \(f\ge1\),
the latent-observation union \(\Callobs{f}\) is not TxtBc-identifiable from
positive data.
\end{theorem}

\begin{proof}
By Corollary~\ref{cor:regular-in-latent-union}, the class contains
\[
  R_k:=\{a,a^2,\ldots,a^k\}\qquad(k\ge1)
\]
and their union \(R_\infty=a^+\).  Since
\(R_1\subsetneq R_2\subsetneq\cdots\),
Lemma~\ref{lem:ascending-chain-obstruction} applies.
\end{proof}

The present section uses observation size as a bound on latent learning advice.
Section~\ref{sec:comparison} returns to the same parameter for a
construction-specific structural calibration: the tagged Rambow--Satta
witnesses used there require observation images that grow with their parameter.
A systematic theory of minimum observation size and its algebraic structure is
outside the scope of the present paper.

\section{Relation to Multidimensional Distributional Learning}
\label{sec:comparison}

Yoshinaka's multidimensional substitutability gives the closest existing
positive-data learning framework for MCFGs.  The role of this section is
calibration: we show that his entire semantic family lies in one
alphabet-dependent observation slice, record explicitly the adaptive-rank
baseline already implicit in his original conjecture construction, and then
calibrate minimum rule rank and observation cost on a tagged Rambow--Satta
family.  We finally compare finite-state observation with finite prefix/suffix
guarding and relate the output-typed proof structure to congruential MCFGs.

\subsection{Yoshinaka's Semantic Family in One Observation Slice}

\begin{definition}[Yoshinaka multicontexts and filling]
\label{def:yoshinaka-multicontexts}
For $m\ge1$, let $\mathsf{Ctx}_{\mathrm Y}^{(m)}$ be the set of expressions
\[
  x=u_0\blank_1u_1\blank_2\cdots
    u_{m-1}\blank_m u_m,
\]
where $u_0,u_m\in\Sigma^*$ and
$u_1,\ldots,u_{m-1}\in\Sigma^+$.  Thus the two outer segments may be empty,
whereas the material between consecutive holes is nonempty.  For a filler
$\tuple y=(y_1,\ldots,y_m)\in(\Sigma^+)^m$, define
\[
  x\odot\tuple y
  :=u_0y_1u_1y_2\cdots u_{m-1}y_mu_m.
\]
This is the nonempty filler/inter-hole convention used in Yoshinaka's journal
definition \cite[Section~3.1, p.~1824]{Yoshinaka2011}.
\end{definition}

\begin{definition}[Yoshinaka's \(p\)D-substitutability
{\cite[Section~3.1, p.~1824]{Yoshinaka2011}}]
\label{def:yoshinaka-pd-substitutability}
A language \(L\) is \emph{\(p\)D-substitutable} if, for every
\(1\le m\le p\), all \(x_1,x_2\in\mathsf{Ctx}_{\mathrm Y}^{(m)}\), and all
\(\tuple y_1,\tuple y_2\in(\Sigma^+)^m\),
\[
 x_1\odot\tuple y_1,\quad x_1\odot\tuple y_2,\quad
 x_2\odot\tuple y_1\in L
 \quad\Longrightarrow\quad
 x_2\odot\tuple y_2\in L.
\]
Write \(S(p)\) for this semantic class.  Yoshinaka's learned class is
\(SL(p,r)=S(p)\cap L(p,r)\), where \(p\) bounds nonterminal dimension and
\(r\) bounds rule-function rank.
\end{definition}

Following Yoshinaka's journal definition, the condition is required for every
arity \(m\le p\), not only for \(m=p\).  Yoshinaka also notes that replacing the
nonempty filler/inter-hole convention by an unrestricted empty-component
version is possible but makes the semantic restriction substantially stronger
\cite[footnote~3, p.~1824]{Yoshinaka2011}.  The boundary construction below
resolves this tension differently: it retains Yoshinaka's original nonempty
\(p\)D class while embedding it into the more permissive named-context
interface used by the reconstruction theorem.

The only superficial mismatch between these multicontexts and one fixed
orientation sector of our named contexts is therefore that Yoshinaka requires
nonempty fillers and nonempty material between consecutive holes.  A finite
boundary quotient removes exactly this mismatch.


\begin{definition}[Finite \((k,\ell)\)-boundary observations]
\label{def:boundary-observation}
Fix integers \(k,\ell\ge0\), put \(s=k+\ell\), and use the conventions
\(\operatorname{pref}_0(w)=\operatorname{suf}_0(w)=\eps\).
For a finite alphabet \(\Sigma\), define an equivalence relation
\(\equiv_{\partial}^{k,\ell}\) on \(\Sigma^*\) by declaring
\(u\equiv_{\partial}^{k,\ell}v\) when either
\begin{enumerate}[label=(\roman*),leftmargin=*]
\item \(u=v\) and \(|u|\le s\), or
\item \(|u|,|v|>s\),
\(\operatorname{pref}_k(u)=\operatorname{pref}_k(v)\), and
\(\operatorname{suf}_\ell(u)=\operatorname{suf}_\ell(v)\).
\end{enumerate}
For \(\Sigma=\emptyset\), use the one-element relation.
\end{definition}

\begin{lemma}[Boundary equivalence is a finite monoid congruence]
\label{lem:boundary-congruence}
For every finite alphabet $\Sigma$ and all $k,\ell\ge0$, the relation
$\equiv_{\partial}^{k,\ell}$ is a finite-index two-sided congruence on the
free monoid $\Sigma^*$.  Hence
\[
  B_\Sigma^{k,\ell}:=\Sigma^*/{\equiv_{\partial}^{k,\ell}}
\]
is a finite monoid and the quotient map
\[
  b_\Sigma^{k,\ell}:\Sigma^*\twoheadrightarrow B_\Sigma^{k,\ell}
\]
is a monoid homomorphism.
\end{lemma}

\begin{proof}
If $u\equiv_{\partial}^{k,\ell}v$ lies in the short part, then $u=v$, so
$xuy=xvy$ for all $x,y\in\Sigma^*$.  Otherwise $|u|,|v|>k+\ell$ and $u,v$
have the same $k$-prefix and the same $\ell$-suffix.  The words $xuy$ and
$xvy$ are again longer than $k+\ell$.  Their first $k$ symbols are either
entirely contained in $x$, or consist of $x$ followed by an initial segment
of length at most $k$ of $u$ or $v$; these agree because the $k$-prefixes of
$u,v$ agree.  The symmetric argument at the right end gives equal
$\ell$-suffixes.  Thus $xuy\equiv_{\partial}^{k,\ell}xvy$.  Finite index
follows because the short part contains only the finitely many words of
length at most $k+\ell$, while every long class is determined by a pair
consisting of a $k$-prefix and an $\ell$-suffix.  The empty-alphabet case is
immediate.
\end{proof}

The boundary observation used in the multidimensional comparison below is the
special case
\[
  b_\Sigma:=b_\Sigma^{1,1}.
\]
Thus \(b_\Sigma\) records words of length at most two exactly and, for longer
words, records only their first and last letters.


\begin{theorem}[Yoshinaka's semantic class is contained in one boundary slice]
\label{thm:yoshinaka-inclusion}
For every finite alphabet \(\Sigma\) and every \(p\ge1\),
\[
  S(p)\cap p\text{-}\mathrm{MCFL}
  \subseteq
  \Cmcf{p}{b_\Sigma}.
\]
Consequently, for every \(r\),
\[
  SL(p,r)\subseteq\Cmcf{p}{b_\Sigma}.
\]
\end{theorem}

\begin{proof}
Let \(L\in S(p)\cap p\text{-}\mathrm{MCFL}\).  Fix \(d\le p\), an
orientation \(\sigma\in S_d\), and tuples \(\tuple x,\tuple y\) with equal
componentwise \(b_\Sigma\)-type that share an accepting context \(E\) of
orientation \(\sigma\).  We show equality of their distributions inside this
same sector.

For each coordinate \(i\), equality of boundary type gives one of two cases.
If \(|x_i|\le2\), then \(x_i=y_i\) exactly.  If \(|x_i|\ge3\), then also
\(|y_i|\ge3\), and there are letters \(a_i,b_i\in\Sigma\) and nonempty
interiors \(x_i^\circ,y_i^\circ\in\Sigma^+\) such that
\[
  x_i=a_i x_i^\circ b_i,
  \qquad
  y_i=a_i y_i^\circ b_i.
\]

Take any context \(F\) of orientation \(\sigma\).  Normalize \(F\) as
follows.  Replace every short-coordinate hole by the common literal word
\(x_i=y_i\).  Replace every remaining long-coordinate hole \(\blank_i\) by
\(a_i\blank b_i\), and then forget the hole names.  Read the surviving holes
from left to right, i.e. in the order inherited from \(\sigma\).  If no holes
survive, then \(F[\tuple x]=F[\tuple y]\) and there is nothing to prove.  If
\(m\ge1\) holes survive, the normalized object belongs to
\(\mathsf{Ctx}_{\mathrm Y}^{(m)}\): every surviving filler is nonempty, and
between two consecutive holes there is at least the last boundary letter of
the earlier long component and the first boundary letter of the later one.
Moreover filling it with the ordered interior tuple obtained from \(\tuple x\)
(resp. \(\tuple y\)) gives exactly \(F[\tuple x]\) (resp. \(F[\tuple y]\)).

Apply the same normalization to the common accepting context \(E\).  Its two
fillings are in \(L\).  If \(F[\tuple x]\in L\), Yoshinaka's
\(m\)D-substitutability, applied to the normalized contexts of \(E\) and
\(F\), gives \(F[\tuple y]\in L\).  Swapping \(\tuple x\) and \(\tuple y\)
gives the converse.  Thus
\[
  \D_{L,\sigma}^{(d)}(\tuple x)
  =\D_{L,\sigma}^{(d)}(\tuple y),
\]
which is Definition~\ref{def:tuple-substitutable}.  Since \(L\) is already a
\(p\)-MCFL, it belongs to \(\Cmcf{p}{b_\Sigma}\).
\end{proof}

\begin{lemma}[Boundary types control the block-order envelope]
\label{lem:block-order-boundary-control}
Let $\Gamma=\{a,b,c\}$ and $P=a^*b^*c^*$.  If $u,v$ are factors of words in
$P$ and $b_\Gamma(u)=b_\Gamma(v)$, then for all $\alpha,\beta\in\Gamma^*$,
\[
  \alpha u\beta\in P
  \quad\Longleftrightarrow\quad
  \alpha v\beta\in P.
\]
\end{lemma}

\begin{proof}
Every factor of a word in $P$ again belongs to $P$.  If $|u|\le2$, equality
of $b_\Gamma$-type gives $u=v$.  Otherwise both words have length at least
three and have the same first and last letters.  A nonempty word of $P$ moves
monotonically through the ordered phases $a<b<c$; once its own block order is
known to be legal, compatibility with arbitrary material on the left and
right depends only on the phase of its first and last letters.  Hence $u$ and
$v$ have the same two-sided behavior with respect to $P$.  The empty-word
case is contained in the exact short part.
\end{proof}

\begin{theorem}[Substantive strictness of the canonical boundary slice]
\label{thm:canonical-boundary-strictness}
Let $\Gamma=\{a,b,c\}$ and $b_\Gamma=b_\Gamma^{1,1}$.  For every integer
$t\ge1$, put
\[
  H_t:=\{a^{tn}b^nc^n:n\ge0\}.
\]
Then every $H_t$ is non-context-free, has minimum fan-out two and minimum
fan-out-two rule rank one, and for every $p\ge2$,
\[
  H_t\in \Cmcf{p}{b_\Gamma}\setminus S(p).
\]
The languages $H_t$ are pairwise distinct.  Consequently, for every $p\ge2$,
\[
  S(p)\cap p\text{-}\mathrm{MCFL}
  \subsetneq
  \Cmcf{p}{b_\Gamma},
\]
and the difference contains an infinite pairwise distinct family of
non-context-free fan-out-two languages of minimum rule rank one.  More
generally, for every nonempty finite alphabet $\Sigma$ and every $p\ge1$,
\[
  S(p)\cap p\text{-}\mathrm{MCFL}
  \subsetneq
  \Cmcf{p}{b_\Sigma}.
\]
\end{theorem}

\begin{proof}
A fan-out-two rank-one grammar for $H_t$ uses a fan-out-two nonterminal $A$
and the rules
\[
  A\to(\eps,\eps),\qquad
  A\to(a^t x_1^1,bx_1^2c)(A),\qquad
  S\to(x_1^1x_1^2)(A).
\]
Thus $H_t\in L(2,1)$ and $\operatorname{rr}_2(H_t)=1$.  It is not
context-free: for the homomorphism
$g:\{u,b,c\}^*\to\Gamma^*$ with
$g(u)=a^t$, $g(b)=b$, and $g(c)=c$,
\[
  g^{-1}(H_t)\cap u^*b^*c^*=\{u^nb^nc^n:n\ge0\},
\]
contradicting closure of CFLs under inverse homomorphism and regular
intersection.  Since fan-out one is exactly the context-free case, the
minimum fan-out is two.

We next prove $(p,b_\Gamma)$-tuple-substitutability, in fact for arbitrary
arity bound.  Put $P=a^*b^*c^*$ and define the additive homomorphism
\[
  \varphi_t(w):=(|w|_a-t|w|_b,\ |w|_b-|w|_c)\in\mathbb Z^2.
\]
Then
\[
  H_t=P\cap\varphi_t^{-1}(0,0).
\]
Fix an arity $d$, an orientation $\sigma$, and tuples
$\tuple x,\tuple y$ of equal componentwise $b_\Gamma$-type sharing an
accepting context $E$ in sector $\sigma$.  Since both $E[\tuple x]$ and
$E[\tuple y]$ lie in $P$, every component $x_i,y_i$ is a factor of a word in
$P$.  Lemma~\ref{lem:block-order-boundary-control}, applied one component at a
time, therefore gives for every context $F$ in the same sector
\[
  F[\tuple x]\in P\quad\Longleftrightarrow\quad F[\tuple y]\in P.
\]
Moreover $E[\tuple x],E[\tuple y]\in\ker\varphi_t$.  Because $\mathbb Z^2$
is abelian, cancellation of the fixed terminal material of $E$ gives
\[
  \sum_{i=1}^d\varphi_t(x_i)=
  \sum_{i=1}^d\varphi_t(y_i).
\]
The same equality implies
$\varphi_t(F[\tuple x])=\varphi_t(F[\tuple y])$ for every $F$, independently
of the orientation.  Hence
\[
  F[\tuple x]\in H_t\quad\Longleftrightarrow\quad F[\tuple y]\in H_t,
\]
so the two sector distributions are equal.  Therefore
$H_t\in\Cmcf{p}{b_\Gamma}$ for every $p\ge2$.

To exclude $S(p)$ it is enough to use arity one.  Let
\[
  y_1=b,\qquad y_2=a^tb^2c,
  \qquad E=a^t\blank_1c,
  \qquad F=a^{2t}\blank_1bc^2.
\]
Then
\[
  E\odot y_1=a^tbc,\qquad
  E\odot y_2=F\odot y_1=a^{2t}b^2c^2
\]
belong to $H_t$, whereas
\[
  F\odot y_2=a^{3t}b^2cbc^2\notin a^*b^*c^*,
\]
so $H_t\notin S(1)$ and hence $H_t\notin S(p)$ for every $p\ge1$.
Finally $H_s\ne H_t$ for $s\ne t$, since $a^tbc\in H_t$ and a word of
$H_s$ with exactly one $b$ must have exactly $s$ occurrences of $a$.
The inclusions themselves are Theorem~\ref{thm:yoshinaka-inclusion}.  For an
arbitrary nonempty $\Sigma$, choose $a\in\Sigma$.  The language $\{a,aa\}$
belongs to $\Cmcf{p}{b_\Sigma}$ because the two words have distinct exact
short types, while the arity-one contexts $E=\blank_1$ and
$F=a\blank_1$ witness failure of $S(1)$, hence of every $S(p)$.
\end{proof}

\paragraph{Boundary-depth calibration.}
The proof of Theorem~\ref{thm:yoshinaka-inclusion} is not special to the
\((1,1)\) boundary.  For every $k,\ell\ge0$ with $k+\ell\ge1$, the same
normalization, using the common $k$-prefix and $\ell$-suffix around every
surviving hole, gives
\[
  S(p)\cap p\text{-}\mathrm{MCFL}
  \subseteq \Cmcf{p}{b_\Sigma^{k,\ell}}.
\]
The condition $k+\ell\ge1$ ensures that material contributed by the normalized
components separates consecutive surviving holes.  Thus every nontrivial
one- or two-sided finite boundary guard already supports the semantic
inclusion; $b_\Sigma=b_\Sigma^{1,1}$ is used below because it is the simplest
symmetric alphabet-only choice.

\begin{remark}[Adaptive rank in Yoshinaka's conjecture construction]
\label{rem:yoshinaka-adaptive-rank}
The qualitative learnability of the union of Yoshinaka's fixed-rank classes is
not specific to the finite-observation construction.  If $\eps\in L$, the empty
word may be handled separately by a start $\eps$-rule once it has appeared; the
sampled tuple nonterminals in Yoshinaka's construction concern nonempty fillers.
Let $Y_{p,q}(K)$ denote
Yoshinaka's conjecture grammar at dimension bound $p$ and rule-rank cap $q$.
Yoshinaka's soundness argument is uniform in $q$; an explicit restatement is
given in \cite[Proposition~1]{CaseYoshinakaZeugmann2013}.  Thus, for every
$p$D-substitutable target $L$ and every finite $K\subseteq L$,
\[
  L(Y_{p,q}(K))\subseteq L.
\]
Moreover, if a target presentation has rule rank $q_*$, Yoshinaka's finite
characteristic-sample argument is complete for every cap $q\ge q_*$
\cite{Yoshinaka2011,CaseYoshinakaZeugmann2013}.  Hence a
single set-driven learner can let the cap grow with its finite content, for
example $q(K)=\max\{1,|K|\}$.  On an infinite target, enlarge the target-side
characteristic sample by finitely many distinct positive examples until its
cardinality is at least $q_*$.  On a finite target, use a lexical presentation
of rank at most one and take the whole finite language if necessary.  After
that finite content has appeared, soundness and completeness give exact
reconstruction.  A conservative wrapper then yields TxtEx convergence by the
same argument as Corollary~\ref{cor:gold-explanatory}.

Thus target-specific rank advice is not, by itself, the novelty boundary with
Yoshinaka.  The point of Theorem~\ref{thm:yoshinaka-inclusion} is instead the
semantic embedding: his whole family is captured by one fixed,
alphabet-dependent observation slice to which the more general reconstruction
theorem applies.  This adaptive-rank baseline is purely qualitative: allowing
$q(K)$ to grow with the sample does not preserve Yoshinaka's fixed-$(p,q)$
polynomial-time guarantee.
\end{remark}

\begin{corollary}[Recovery of Yoshinaka's semantic family]
\label{cor:yoshinaka-full-union-learning}
For every fixed finite alphabet $\Sigma$ and fan-out bound $p$, the canonical
finite-observation learner supplied with $b_\Sigma$ TxtBc-identifies
\[
  S(p)\cap p\text{-}\mathrm{MCFL}
  =\bigcup_{q\ge1}SL(p,q)
\]
from positive data.  This is recorded as a calibration consequence, not as a
new union-learnability result.
\end{corollary}

\begin{proof}
By Theorem~\ref{thm:yoshinaka-inclusion}, every target in the displayed family
belongs to the single slice $\Cmcf{p}{b_\Sigma}$.  Corollary~\ref{cor:semantic-identification}
therefore applies uniformly to the whole family.
\end{proof}

\subsection{Rule Rank and Observation Cost: Tagged and Fixed-Alphabet Calibrations}
\label{subsec:rank-observation-calibration}

The next construction calibrates the Yoshinaka embedding against minimum rule
rank.  We first use the tagged Rambow--Satta family because its exact rank
hierarchy is transparent, and then compress that hierarchy to one fixed
three-letter alphabet.  The latter shows that alphabet growth is not the
source of unbounded minimum rank.  We then quantify the genuinely remaining
obstruction: the finite observation needed by this construction must still
grow with the parameter.

For $r\ge2$ and $\pi\in S_r$, let
\[
 P_{r,\pi}=c_1^*d_1c_2^*d_2\cdots d_{2r-1}c_{2r}^*,
\]
where
$c_1,\ldots,c_{2r}=A_1,\ldots,A_r,B_{\pi(1)},\ldots,B_{\pi(r)}$
and all delimiters $d_j$ are distinct.  Let
$\delta(A_i)=e_i$, $\delta(B_i)=-e_i$, and $\delta(d_j)=0$ in $\mathbb Z^r$,
and put $T_{r,\pi}=P_{r,\pi}\cap\ker\delta$.

\begin{theorem}[Strict rule-rank hierarchy inside $S(2)$]
\label{thm:s2-rank-hierarchy}
For every $r\ge6$ and the Rambow--Satta hierarchy permutation $\pi_r$,
\[
  T_{r,\pi_r}\in L(2,r-2)\setminus L(2,r-3).
\]
Moreover $T_{r,\pi_r}\in S(2)$, and hence
\[
  T_{r,\pi_r}\in SL(2,r-2)\setminus SL(2,r-3).
\]
Consequently minimum rule rank is unbounded on
$S(2)\cap2\text{-}\mathrm{MCFL}$ across this alphabet-indexed family.
\end{theorem}

\begin{proof}[Proof sketch]
The unique-delimiter construction places $P_{r,\pi_r}$ in $S(2)$, and
intersecting with the zero fiber of the abelian morphism $\delta$ preserves
multidimensional substitutability, giving $T_{r,\pi_r}\in S(2)$.  Erasing the
delimiters maps $T_{r,\pi_r}$ onto the corresponding zero-exponent
Rambow--Satta witness; closure under homomorphism within the same fixed
fan-out/rule-rank stratum gives the lower bound, while inverse homomorphism and
regular intersection within that stratum transfer the matching upper bound.  Appendix~\ref{app:rank-hierarchy} gives the details.
\end{proof}


\begin{corollary}[The rank hierarchy persists over one fixed alphabet]
\label{cor:fixed-alphabet-rank-hierarchy}
Let $\Delta=\{a,b,c\}$.  For $r\ge6$ and the Rambow--Satta hierarchy
permutation $\pi_r$, define
\[
 C_r:=\{a^{n_1}c\cdots ca^{n_r}c
 b^{n_{\pi_r(1)}}c\cdots cb^{n_{\pi_r(r)}}:
 n_1,\ldots,n_r\ge0\}.
\]
Then
\[
  C_r\in L(2,r-2)\setminus L(2,r-3).
\]
Hence minimum rule rank is unbounded already over the fixed alphabet $\Delta$.
\end{corollary}

\begin{proof}
Let $q_r$ map every $A_i$ to $a$, every $B_i$ to $b$, and every delimiter
$d_j$ to $c$.  Then $q_r(T_{r,\pi_r})=C_r$, so the homomorphism closure of
the fixed fan-out/rank stratum and
Theorem~\ref{thm:s2-rank-hierarchy} give $C_r\in L(2,r-2)$.

Suppose $C_r\in L(2,r-3)$.  On the source alphabet,
\[
  T_{r,\pi_r}=q_r^{-1}(C_r)\cap P_{r,\pi_r}.
\]
The fixed-parameter stratum $L(2,r-3)$ is closed under inverse homomorphism and
intersection with regular languages by the Rambow--Satta full-AFL theorem and
its MCFG correspondence; the same closure facts are recalled in
Appendix~\ref{app:rank-hierarchy}.  Hence
$T_{r,\pi_r}\in L(2,r-3)$, contradicting
Theorem~\ref{thm:s2-rank-hierarchy}.
\end{proof}

The preceding recognizable-language proposition gives a general sufficient
condition for membership in a fixed-observation slice.  The next proposition
gives a complementary necessary capacity bound: one finite observation can
support only finitely many pairwise distribution-distinct tuple types inside a
single shared context.

\begin{proposition}[Shared-context type-packing bound]
\label{prop:shared-context-packing}
Let $L$ be orientation-aware $(f,h)$-tuple-substitutable, where
$h:\Sigma^*\to M$ is finite.  Fix $1\le d\le f$ and an orientation
$\sigma\in S_d$.  Suppose that
\[
  \tuple x_1,\ldots,\tuple x_N\in(\Sigma^*)^d
\]
share one accepting context
\[
  E\in\bigcap_{i=1}^N \D_{L,\sigma}^{(d)}(\tuple x_i)
\]
and have pairwise distinct sector distributions:
\[
  \D_{L,\sigma}^{(d)}(\tuple x_i)
  \ne
  \D_{L,\sigma}^{(d)}(\tuple x_j)
  \qquad(i\ne j).
\]
Then
\[
  N\le |\operatorname{im}h|^d.
\]
\end{proposition}

\begin{proof}
If
\[
  h^{(d)}(\tuple x_i)=h^{(d)}(\tuple x_j)
\]
for some $i\ne j$, then the common accepting context $E$ makes the two sector
distributions intersect.  Orientation-aware $(f,h)$-tuple-substitutability
would therefore force
\[
  \D_{L,\sigma}^{(d)}(\tuple x_i)
  =
  \D_{L,\sigma}^{(d)}(\tuple x_j),
\]
contrary to the hypothesis.  Hence the $N$ tuples have pairwise distinct
$h^{(d)}$-types.  There are at most $|\operatorname{im}h|^d$ such types.
\end{proof}




\begin{corollary}[Fixed-alphabet persistence of the observation obstruction]
\label{cor:fixed-alphabet-rank-observation}
For the fixed-alphabet languages $C_r$ of
Corollary~\ref{cor:fixed-alphabet-rank-hierarchy}, every finite observation
$h:\Delta^*\to M$ with $C_r\in\Cmcf{2}{h}$ satisfies
\[
  r\le |\operatorname{im}h|^2.
\]
Thus the exact minimum-rank hierarchy survives alphabet compression, but no
single finite observation witnesses this particular family for unbounded
$r$.
\end{corollary}

\begin{proof}
Put $D=c^{2r-1}$ and, for $1\le j\le r$, let
\[
  \tau_j=(c^{j-1},c^{r-j}).
\]
All $\tau_j$ share the identity-oriented accepting context
$E=\blank_1\blank_2c^r$, since $E[\tau_j]=D\in C_r$.
For $s(j)=\pi_r^{-1}(j)$ define
\[
  F_j:=\blank_1a\blank_2c^{s(j)}bc^{r-s(j)}.
\]
Then $F_j[\tau_j]\in C_r$: the unique nonzero first-half exponent is the
$j$th one and the unique nonzero second-half exponent is its matching
$\pi_r$-position.  If $k\ne j$, then $F_j[\tau_k]\notin C_r$, since the
nonzero first-half exponent moves to position $k$ while the nonzero
second-half exponent still corresponds to $j$.  Hence the $r$ tuples have a
common accepting context and pairwise distinct identity-sector distributions.
Proposition~\ref{prop:shared-context-packing} with $d=2$ gives the claim.
\end{proof}

Corollaries~\ref{cor:fixed-alphabet-rank-hierarchy}
and~\ref{cor:fixed-alphabet-rank-observation} sharpen the interpretation:
the rank hierarchy is not an artifact of a growing alphabet, yet the
observation obstruction persists after compression to the fixed alphabet
$\{a,b,c\}$.  This is a construction-specific fact and should not be read as
a general theorem that high minimum rule rank forces a large observation, nor
as a bound on minimum rank for arbitrary members of $\Clatent{2}{k}$.
Whether one fixed finite-observation slice at fan-out two can have unbounded
minimum rule rank remains open.

Thus no one finite rule-rank cap covers this calibration even on a fixed
three-letter alphabet, but this particular hierarchy also cannot be witnessed
by one uniformly bounded finite observation.  In view of
Remark~\ref{rem:yoshinaka-adaptive-rank}, it is used as structural calibration,
not as evidence that adaptive removal of Yoshinaka's rank parameter is itself
a new learnability result.

\subsection{Finite-State Control versus Finite Boundary Guards}

To compare the general finite-state interface with bounded prefix/suffix
information, we use the following direct multidimensional guarded analogue of
Yoshinaka's $(k,\ell)$ condition.  It is introduced only as a controlled
comparison class; it is not attributed to Yoshinaka.

\begin{definition}[Multidimensional \((k,\ell)\)-guarded substitutability]
\label{def:guarded-pd-substitutability}
Fix \(p\ge1\) and \(k,\ell\ge0\).  A language \(L\subseteq\Sigma^*\) is
\emph{\((k,\ell)\)-guarded \(p\)D-substitutable} if, for every
\(1\le m\le p\), all \(x_1,x_2\in\mathsf{Ctx}_{\mathrm Y}^{(m)}\), and all
\(\tuple y_1,\tuple y_2\in(\Sigma^+)^m\) satisfying
\[
  b_\Sigma^{k,\ell}(y_{1,i})=b_\Sigma^{k,\ell}(y_{2,i})
  \qquad(1\le i\le m),
\]
the implication
\[
 x_1\odot\tuple y_1,\quad x_1\odot\tuple y_2,\quad
 x_2\odot\tuple y_1\in L
 \quad\Longrightarrow\quad
 x_2\odot\tuple y_2\in L
\]
holds.  Write \(S^{k,\ell}(p)\) for this semantic class.  This is introduced
as the direct multidimensional analogue of Yoshinaka's
\((k,\ell)\)-substitutability \cite{Yoshinaka2008}, not as a claim that this
exact multidimensional variant was already isolated in that work.  Since all
fillers are nonempty, \(S^{0,0}(p)=S(p)\).
\end{definition}


\begin{proposition}[Same-depth boundary slices lie inside the guarded class]
\label{prop:same-depth-guarded-converse}
For every finite alphabet $\Sigma$, every $p\ge1$, and all $k,\ell\ge0$,
\[
  \Cmcf{p}{b_\Sigma^{k,\ell}}
  \subseteq
  S^{k,\ell}(p)\cap p\text{-}\mathrm{MCFL}.
\]
\end{proposition}

\begin{proof}
Let $L\in\Cmcf{p}{b_\Sigma^{k,\ell}}$ and suppose the premises of
Definition~\ref{def:guarded-pd-substitutability} hold at arity $m\le p$ for
multicontexts $x_1,x_2$ and positive filler tuples
$\tuple y_1,\tuple y_2$ of equal componentwise boundary type.  The
multicontext $x_1$ is an identity-oriented named context accepting both
fillers, so their identity-sector distributions intersect.  Orientation-aware
tuple substitutability therefore makes those distributions equal.  Since
$x_2$ accepts $\tuple y_1$, it also accepts $\tuple y_2$.  Thus the guarded
rectangular implication holds.  The language is already a $p$-MCFL by the
definition of the slice.
\end{proof}


\begin{theorem}[Buffered embedding of guarded multidimensional substitutability]
\label{thm:guarded-buffered-embedding}
For every finite alphabet \(\Sigma\), every \(p\ge1\), and all
\(k,\ell\ge0\),
\[
  S^{k,\ell}(p)\cap p\text{-}\mathrm{MCFL}
  \subseteq
  \Cmcf{p}{b_\Sigma^{k,\ell+1}}.
\]
By left--right symmetry one also has
\[
  S^{k,\ell}(p)\cap p\text{-}\mathrm{MCFL}
  \subseteq
  \Cmcf{p}{b_\Sigma^{k+1,\ell}}.
\]
Thus one additional boundary symbol, on either side, is sufficient to convert
the named-context interface into the nonempty-inter-hole interface of
Yoshinaka multicontexts.
\end{theorem}

\begin{proof}
We prove the suffix-buffered statement; the prefix-buffered statement is its
mirror image.  Fix \(1\le d\le p\), an orientation \(\sigma\in S_d\), and
tuples \(\tuple x,\tuple y\in(\Sigma^*)^d\) of equal componentwise
\(b_\Sigma^{k,\ell+1}\)-type that share an accepting named context
\(E\in\mathcal E_{d,\sigma}\).  Let \(F\in\mathcal E_{d,\sigma}\) and assume
\(F[\tuple x]\in L\); it suffices to prove \(F[\tuple y]\in L\).

If a component \(x_i\) lies in the short exact part of the boundary quotient,
then \(x_i=y_i\); replace that hole in both \(E\) and \(F\) by this common
literal.  For every surviving component, both words are longer than
\(k+\ell+1\).  Equality of boundary type therefore gives a common final
letter \(c_i\) and factorizations
\[
  x_i=x_i'c_i,\qquad y_i=y_i'c_i,
\]
with \(x_i',y_i'\ne\eps\) and
\[
  b_\Sigma^{k,\ell}(x_i')=b_\Sigma^{k,\ell}(y_i').
\]
Replace the corresponding named hole by \(\blank_i c_i\).  In surface order,
every pair of consecutive surviving holes is now separated by at least the
buffer letter contributed by the earlier hole.  Let \(m\le d\) be the number of surviving holes.  If \(m=0\), then all
components were in the exact short part, so \(\tuple x=\tuple y\) on every
hole and the claim is immediate.  Suppose \(m\ge1\).  After forgetting hole
names, the normalized versions of \(E\) and \(F\) are members of
\(\mathsf{Ctx}_{\mathrm Y}^{(m)}\).  Read the surviving components in their
surface order, equivalently in the order induced by the common orientation
\(\sigma\); the resulting filler tuples are nonempty and have equal
componentwise \((k,\ell)\)-boundary types.  The three accepted fillings
\(E[\tuple x],E[\tuple y],F[\tuple x]\) now satisfy
Definition~\ref{def:guarded-pd-substitutability}, hence
\(F[\tuple y]\in L\).  Interchanging \(\tuple x\) and \(\tuple y\) gives the
converse, so the orientation-sector distributions are equal.  Since
\(d,\sigma\) were arbitrary and \(L\) is a \(p\)-MCFL, the claim follows.
\end{proof}

Combining Proposition~\ref{prop:same-depth-guarded-converse} with
Theorem~\ref{thm:guarded-buffered-embedding} gives the same-depth/buffered
sandwich
\[
  \Cmcf{p}{b_\Sigma^{k,\ell}}
  \subseteq
  S^{k,\ell}(p)\cap p\text{-}\mathrm{MCFL}
  \subseteq
  \Cmcf{p}{b_\Sigma^{k,\ell+1}},
\]
with the symmetric final term $b_\Sigma^{k+1,\ell}$ also available.


\begin{theorem}[An infinite non-context-free separation from all finite boundary guards]
\label{thm:guarded-infinite-noncfl-separation}
Let \(\Gamma=\{a,b,c\}\) and let
\[
 R_\star:=\{a^i b^j c^k:i,k\ge0,\ j\ge1,\ j\bmod3\in\{0,1\}\}.
\]
Fix once and for all the transition morphism
\(h_\star:\Gamma^*\to M_\star\) of a finite automaton recognizing
\(R_\star\).  For every integer \(t\ge1\), put
\[
  L_t:=\{a^{tn}b^nc^n:n\ge1,\ n\bmod3\in\{0,1\}\}.
\]
Then the following hold.
\begin{enumerate}[label=(\roman*),leftmargin=*]
\item Every \(L_t\) is a non-context-free \(2\)-MCFL with a fan-out-two
rule-rank-one presentation.
\item Every \(L_t\) belongs to the same fixed observation slice
\(\Cmcf{p}{h_\star}\) for every \(p\ge2\).
\item For every finite \(k,\ell\ge0\) and every \(p\ge2\),
\(L_t\notin S^{k,\ell}(p)\).
\end{enumerate}
Consequently, for every \(p\ge2\), \(q\ge1\), and \(k,\ell\ge0\), if
\[
 \widehat h_{k,\ell}:=(b_\Gamma^{k,\ell+1},h_\star),
\]
where the pair denotes the product homomorphism
\[
  w\longmapsto\bigl(b_\Gamma^{k,\ell+1}(w),h_\star(w)\bigr)
  \quad\text{into}\quad B_\Gamma^{k,\ell+1}\times M_\star,
\]
then
\[
 S^{k,\ell}(p)\cap L(p,q)
 \subsetneq
 \Cmcf{p}{\widehat h_{k,\ell}}\cap L(p,q),
\]
and the difference contains the infinite pairwise distinct family
\(\{L_t:t\ge1\}\) of non-context-free rank-one languages.
\end{theorem}

\begin{proof}
A rank-one fan-out-two presentation for \(L_t\) uses a fan-out-two
nonterminal \(A\), the two rank-zero rules
\[
 A\to(a^t,bc),\qquad A\to(a^{3t},b^3c^3),
\]
the rank-one recursive rule
\[
 A\to(a^{3t}x_1^1,b^3x_1^2c^3)(A),
\]
and the start rule \(S\to(x_1^1x_1^2)(A)\).  Hence \(L_t\in L(2,1)\).
It is not context-free.  Indeed, under the homomorphism
\(g:\{u,b,c\}^*\to\Gamma^*\) with
\(g(u)=a^t\), \(g(b)=b\), and \(g(c)=c\), inverse homomorphism followed by
intersection with \(u^*b^*c^*\) would make
\[
 \{u^nb^nc^n:n\ge1,\ n\bmod3\in\{0,1\}\}
\]
context-free if \(L_t\) were.  Intersecting further with
\((u^3)^*(b^3)^*(c^3)^*\) and applying the inverse of the length-tripling
homomorphism would then make \(\{u^mb^mc^m:m\ge1\}\) context-free, a
contradiction.

For slice membership define the additive homomorphism
\[
 \varphi_t(w):=(|w|_a-t|w|_b,\ |w|_b-|w|_c):\Gamma^*\to\mathbb Z^2.
\]
Then
\[
 L_t=R_\star\cap\varphi_t^{-1}(\{(0,0)\}).
\]
The fixed observation \(h_\star\) recognizes \(R_\star\), so
Proposition~\ref{prop:regular-abelian-balance} yields
\(L_t\in\Cmcf{p}{h_\star}\) for every \(p\ge2\).  Notice that the same
observation works for all \(t\).  It is nevertheless chosen from the common
regular envelope \(R_\star\) of this separating family, rather than canonically
from the alphabet alone.

It remains to exclude every finite boundary guard.  Fix \(k,\ell\ge0\) and
choose \(n=3m\) so large that all boundary prefixes and suffixes below lie in
the long part of \(b_\Gamma^{k,\ell}\).  Put
\[
 \tuple x=(a^{tn}b^{n-1},c^n),\qquad
 \tuple y=(a^{t(n+1)}b^n,c^{n+1}).
\]
Then \(b_\Gamma^{k,\ell}\) agrees componentwise on \(\tuple x,\tuple y\).
With the two Yoshinaka multicontexts
\[
 E=\blank_1 b\blank_2,
 \qquad
 F=a^{2t}\blank_1 b^3\blank_2c^2,
\]
we have
\[
 E\odot\tuple x=a^{tn}b^nc^n\in L_t,
 \qquad
 E\odot\tuple y=a^{t(n+1)}b^{n+1}c^{n+1}\in L_t,
\]
and
\[
 F\odot\tuple y=a^{t(n+3)}b^{n+3}c^{n+3}\in L_t,
 \qquad
 F\odot\tuple x=a^{t(n+2)}b^{n+2}c^{n+2}\notin L_t.
\]
Thus Definition~\ref{def:guarded-pd-substitutability} fails after exchanging
the names of the two filler tuples in its implication.  Hence
\(L_t\notin S^{k,\ell}(2)\), and therefore not in \(S^{k,\ell}(p)\) for
any \(p\ge2\).

Finally, Theorem~\ref{thm:guarded-buffered-embedding} gives
\(S^{k,\ell}(p)\cap L(p,q)\subseteq
\Cmcf{p}{b_\Gamma^{k,\ell+1}}\cap L(p,q)\).  The product observation
\(\widehat h_{k,\ell}\) refines both factors, and every \(L_t\) belongs to
its slice while lying in \(L(p,1)\subseteq L(p,q)\).  The languages are
pairwise distinct because every word of \(L_t\) satisfies
\(|w|_a=t|w|_b\).  This proves strictness and infinitude of the difference.
\end{proof}


\begin{remark}[Boundary depth versus general finite-state control]
Theorem~\ref{thm:canonical-boundary-strictness} and
Proposition~\ref{prop:same-depth-guarded-converse} separate two phenomena.
The canonical slice \(\Cmcf{p}{b_\Gamma}\) already extends Yoshinaka's
unrestricted semantic family \(S(p)\) by an infinite non-context-free
rank-one family.  But a language in the same-depth slice
\(\Cmcf{p}{b_\Gamma^{k,\ell}}\) automatically belongs to the corresponding
guarded class \(S^{k,\ell}(p)\).  Hence no same-depth boundary observation can
witness separation from that guard.  The fixed morphism \(h_\star\) in
Theorem~\ref{thm:guarded-infinite-noncfl-separation} records an internal
regular property of the middle block and supplies genuinely finite-state
information beyond bounded boundary data; it is used precisely to obtain
separation from every finite guard.
\end{remark}

\subsection{Congruential Decomposition}
\begin{proposition}[Rank-preserving congruential decomposition]
\label{prop:congruential-decomposition}
Let \(L\in\Cmcf{f}{h}\cap L(f,q)\).  Then there are languages
\(L_1,\ldots,L_m\), with
\(m\le |\operatorname{im}h|\), such that
\[
  L=\bigcup_{i=1}^m L_i,
\]
and every \(L_i\) has a congruential MCFG presentation of fan-out at most
\(f\) and rule rank at most \(q\).  Here congruential is meant in the standard
nonpermuting MCFG sense: for every nonterminal, all tuples it generates have
the same ordered language-context set
\cite{Yoshinaka2010,ClarkYoshinaka2016}.
\end{proposition}

\begin{proof}
Start with a fan-out-\(f\), rule-rank-\(q\) presentation of \(L\), apply the
same rank-preserving nondeleting/nonpermuting normalization used in the proof
construction, and form the trimmed output-type refinement
\(\widetilde G_0\) of Definition~\ref{def:output-refinement}.  For every
surviving typed nonterminal \(X=A_{\mathbf p}\),
Proposition~\ref{prop:output-invariants} gives one componentwise \(h\)-type
\(\mathbf p\), while Lemma~\ref{lem:identity-exposing-contexts} supplies a
single identity-oriented context \(\chi(X)\) accepting every tuple in
\(L_X\).  Fixed-\(h\) substitutability therefore makes all tuples in \(L_X\)
have the same ordered context set in \(L\).

Now remove the fresh start symbol \(\widetilde S\) and, for each surviving
typed start state \(S_{(s)}\), take the reachable trimmed subgrammar with
\(S_{(s)}\) as its single start symbol.  Its language is
\[
  L_s=L\cap h^{-1}(s),
\]
by unique lifting of derivations.  The same congruential conclusion holds
relative to \(L_s\): if two tuples have the same typed \(h\)-value, then every
ordered context gives fillings with the same total \(h\)-value, so restricting
from \(L\) to the fiber \(h^{-1}(s)\) preserves equality of their context
sets.  Thus each component grammar is congruential.  Output typing changes no
rule template, tuple dimension, or rule arity, and deleting the fresh start
symbol adds nothing, so fan-out and rule rank remain bounded by \(f\) and
\(q\).  There are at most \(|\operatorname{im}h|\) surviving unary start
types, and their languages union to \(L\).
\end{proof}

\paragraph{Positive bridge to the syntactic approach.}
The proposition gives a presentation-level connection complementary to
Clark--Yoshinaka's language-derived algebraic semantics: the externally
supplied finite observation does not replace the syntactic structure, but
finitely colors and separates congruential components that are exposed by
positive contexts, without increasing fan-out or rule rank.  The whole target
need not itself be a single congruential component under the standard
single-start convention, which is why the finite decomposition is the natural
statement here.  This rank preservation is also distinct from the pcis/SCL
minimality of \cite{ClarkYoshinaka2014Algebraic}, which concerns faithful
semantic structure and nonterminal identification rather than minimum
right-hand-side rank over all equivalent grammars.

\paragraph{Canonicality and strong learning under structural restrictions.}
Clark gives a strong learner for a class of probabilistic dimension-two MCFGs
from positive samples of strings \cite{Clark2021Strong}.  The learnable class
uses well-nested normal forms together with double anchoring, local
unambiguity, and completeness; the admissible production sets supplied to the
algorithm implicitly fix an upper bound on rule rank.  Under these conditions
the learned grammar converges, in the stated probabilistic sense, to a grammar
isomorphic to the target rather than merely to one generating the same string
language.  This result is complementary to the present one.  Clark recovers a
canonical grammar presentation from strings under strong target-side
structural restrictions, whereas our learner identifies the target language
inside an externally supplied finite-observation slice and is not given a
target-specific rank bound; we assume neither well-nestedness nor anchoring or
completeness.  The rank questions are correspondingly different: dimension-two
well-nested MCFGs admit binarization without changing their weak generative
capacity \cite{Clark2021Strong}, while the open problem here concerns minimum
rule rank in unrestricted fixed-observation slices.


\subsection{Other Distributional Learning Frameworks}

The present semantic slices are related to, but not identical with, the
finite-kernel and finite-context approaches to distributional MCFG learning
\cite{Yoshinaka2010,ClarkYoshinaka2014,ClarkYoshinaka2016,KanazawaYoshinaka2017}.
Those conditions are normally imposed on a selected grammar presentation,
whereas $\Cmcf{f}{h}$ is a language-level semantic class determined by an
external finite observation.  The output-type refinement has a
single-representative flavor inside each typed component, but we do not claim a
finite-context characterization of the whole slice.

For a language-level comparison, let $\mathsf{FCP}_{\rm PMCFG}$ denote the
union of Clark--Yoshinaka finite-context PMCFG classes over all finite
parameters \cite{ClarkYoshinaka2012PMCFG,ClarkYoshinaka2014}.

\begin{proposition}[Finite-context PMCFGs are not contained in finite-observation slices]
\label{prop:fcp-not-contained-observation}
For $\Sigma=\{a,b\}$, let $\mathrm{COPY}_*:=\{ww:w\in\Sigma^*\}$.  Then
$\mathrm{COPY}_*\in\mathsf{FCP}_{\rm PMCFG}$ but
\[
 \mathrm{COPY}_*\notin
 \bigcup_{f\ge2}\ \bigcup_{\substack{h:\Sigma^*\to M\\ M\text{ finite}}}
 \Cmcf{f}{h}.
\]
Hence $\mathsf{FCP}_{\rm PMCFG}$ is not contained in the union of finite-observation
MCFG slices.
\end{proposition}

\begin{proof}
Clark and Yoshinaka place the copy language in their finite-context class
$L(1,1,2,1)$ \cite[Definition~2 and Example~3]{ClarkYoshinaka2012PMCFG}.
The witness in Proposition~\ref{prop:copy-infinite-observation} uses only
nonempty words, so it applies unchanged after adjoining $\eps$; thus
$\mathrm{COPY}_*$ fails $(2,h)$-tuple substitutability for every finite $h$.
The same arity-two witness excludes every fan-out bound $f\ge2$.
\end{proof}

Clark's strong-learning result for a probabilistic dimension-two MCFG class
uses well-nested normal forms together with anchoring, local unambiguity, and
completeness, and recovers a canonical grammar presentation
\cite{Clark2021Strong}.  Our guarantee is different: we identify the target
language in a finite-observation slice, without receiving a target-specific
rule-rank bound, but do not claim grammar isomorphism or the same quantitative
data bounds.  Positive-data-plus-query approaches provide another orthogonal
tradeoff by replacing external finite observation with active membership
information \cite{Yoshinaka2010,KanazawaYoshinaka2021,KanazawaYoshinaka2023,YoshinakaClark2012}.

At fan-out one, the same finite-typing idea specializes to the author's
context-free setting \cite{kuriyamaCFG}.  There, standard CFG binarization
removes right-hand-side rank as an independent structural parameter.  This is
not available for general MCFGs of fan-out at least two, where minimum rule
rank forms a genuine hierarchy \cite{RambowSatta1999}.  The
present proof is therefore not merely an invocation of the CFG result: tuple
orientation, induced child orientation, silent-child compression, and
arbitrary-rank MCFG composition are essential to the multidimensional
reconstruction argument developed here.


\section{Conclusion}
\label{sec:conclusion}

Finite observation provides a sufficient compositional interface for
positive-data reconstruction of a broad semantic family of multiple
context-free languages.  For every fixed fan-out and supplied finite
observation, the canonical learner exactly reconstructs the target after a
finite characteristic sample without receiving a target-specific rule-rank
bound.  The proof isolates the mechanism responsible for this removal:
output typing and orientation make semantic substitutions compositional,
silent-child compression leaves only positively supported children, and the
sample itself bounds the rank of every rule witness that must be enumerated.
A conservative wrapper gives TxtEx convergence, while a fixed external
branching cap yields polynomial-time construction of the current hypothesis.

Treating the observation as learning advice yields a complementary boundary.
If the actual observation is hidden but its codomain size is bounded, a single
effective universal refinement restores TxtBc learnability.  This does not
make the information cheap: bounded slices already force
$2^{\Omega(\sqrt{k})}$ positive locking information for arbitrary TxtBc
learners along an explicit family.  If the size bound is removed entirely, the
latent union is not TxtBc-identifiable from positive data.

The comparison with multidimensional distributional learning shows that these
results are not obtained by choosing an artificially narrow semantic slice.
On $\Gamma=\{a,b,c\}$, the canonical symmetric boundary observation
$b_\Gamma=b_\Gamma^{1,1}$ already gives, for every $p\ge2$,
\[
  S(p)\cap p\text{-}\mathrm{MCFL}
  \subsetneq
  \Cmcf{p}{b_\Gamma},
\]
and the difference contains the infinite pairwise distinct family
$H_t=\{a^{tn}b^nc^n:n\ge0\}$ of non-context-free fan-out-two languages of
minimum rule rank one.  The same-depth/buffered sandwich with the guarded
classes then shows exactly why the stronger separation from every finite
boundary guard requires the additional fixed finite-state factor $h_\star$.
The qualitative union learnability already has the adaptive-rank explanation
recorded in Remark~\ref{rem:yoshinaka-adaptive-rank}; the new learning theorem
is instead for the broader fixed-observation slices.

The Rambow--Satta calibration can be compressed to the fixed alphabet
$\{a,b,c\}$ while preserving the exact minimum rule rank
$r-2$.  Thus its unbounded rank is not an alphabet-growth artifact.  However,
the shared-context packing obstruction also survives this compression: every
finite observation witnessing the parameter-$r$ fixed-alphabet member
satisfies $r\le|\operatorname{im}h|^2$.  This remains a construction-specific calibration rather than a general theorem
that high rule rank forces large observation.  In particular, it still does
not establish unbounded minimum rule rank inside one fixed observation slice.  The rank-preserving
congruential decomposition further connects the external observation
interface with language-derived distributional structure.

Several questions remain.  First, natural fixed-observation subclasses may
admit stronger quantitative bounds on the total length of characteristic data
than are proved here.  Second, it is open which weaker forms of finite
compositional advice still suffice for rank-parameter-free reconstruction.
Third, a sharper language-level comparison with finite-kernel and finite-context
MCFG learning may identify further regions where external finite observation
and presentation-derived distributional information can be exchanged.  A
sharper structural question is whether there exist a fixed finite alphabet
$\Sigma$ and a fixed finite observation $h:\Sigma^*\to M$ such that
\[
  \sup\{\operatorname{rr}_2(L):L\in\Cmcf{2}{h}\}=\infty.
\]
The Rambow--Satta calibration does not answer this question even after
fixed-alphabet compression: Corollary~\ref{cor:fixed-alphabet-rank-observation}
forces every observation witnessing its parameter-$r$ member to have image
size at least $\sqrt r$.  Thus alphabet growth can be removed, but the
observation growth cannot be removed by this construction.  It remains open
whether any one fixed slice can support unbounded minimum rule rank.  More
generally, it remains open to characterize how far
individual fixed-observation slices extend beyond the concrete algebraic and
distributional subclasses isolated here, and whether every such slice is
exhausted by combinations of regular envelopes and algebraic fibers.  A
systematic theory of the minimum size and algebraic structure of witnessing
observations is also left open.

\section*{Acknowledgments}
The author thanks Makoto Kanazawa.  He is especially grateful to Ryo Yoshinaka
for encouraging him to return to and complete the doctoral dissertation he had
left unfinished ten years earlier.

\appendix
\section{Technical Proof of the Rank Hierarchy}
\label{app:rank-hierarchy}

This appendix supplies the details used in
Theorem~\ref{thm:s2-rank-hierarchy}.  We separate the two ingredients: the
tagged languages remain multidimensionally substitutable, and tagging preserves
the exact Rambow--Satta minimum rule rank.

For pairwise distinct phase letters $c_1,\ldots,c_k$ and delimiters
$d_1,\ldots,d_{k-1}$, put
\[
  P_k=c_1^*d_1c_2^*d_2\cdots d_{k-1}c_k^*.
\]

\begin{lemma}[Unique-delimiter chain]
\label{lem:unique-delimiter-chain}
For every $p\ge1$, the regular language $P_k$ belongs to $S(p)$.
\end{lemma}

\begin{proof}
Use the partial DFA with states $1,\ldots,k$ and only the transitions
\[
  i\xrightarrow{c_i}i,
  \qquad
  i\xrightarrow{d_i}i+1.
\]
Because every phase letter and delimiter is unique, every nonempty word $u$
induces a partial transition whose domain has size at most one.  In
Yoshinaka's $p$-dimensional multicontexts the fillers and all material between
consecutive holes are nonempty.  Suppose two filling tuples share one
accepting multicontext.  The first hole has the same incoming state in the two
runs because the fixed prefix is the same.  For every later hole, the fixed
nonempty material immediately preceding it has a partial transition with at
most one possible start state, so acceptance forces the same incoming state
in both runs.  Dually, every hole except the last has the same outgoing state
because the following fixed inter-hole material is nonempty and can start in
at most one state; for the last hole the common fixed suffix together with the
unique final state gives the same conclusion.  Thus corresponding fillers
realize the same state-to-state transition.  Replacing them inside any other
accepting multicontext preserves the entire state sequence.  The two tuples
have equal distributions, so $P_k\in S(p)$.
\end{proof}

\begin{lemma}[Abelian-fiber preservation]
\label{lem:abelian-fiber-preservation}
Let $P\in S(p)$ and let $\delta:\Sigma^*\to A$ be a morphism to an abelian
group.  For every $g\in A$,
\[
  P\cap\delta^{-1}(g)\in S(p).
\]
\end{lemma}

\begin{proof}
Suppose two filling tuples $\tuple y,\tuple y'$ share a $p$-multicontext $x$
and both filled words lie in $P\cap\delta^{-1}(g)$.  Since $A$ is abelian,
the fixed material of $x$ contributes the same group element in both words,
and therefore
\[
  \sum_i\delta(y_i)=\sum_i\delta(y_i').
\]
For any other multicontext $x'$, membership of $x'\odot\tuple y$ in $P$
implies membership of $x'\odot\tuple y'$ in $P$ because $P\in S(p)$; the
displayed equality also makes their $\delta$-values equal.  Thus membership in
the fiber $\delta^{-1}(g)$ transfers as well.  The converse is symmetric.
\end{proof}

For $r\ge2$ and $\pi\in S_r$, take the $2r$ phases
\[
  c_1,\ldots,c_{2r}
  =A_1,\ldots,A_r,B_{\pi(1)},\ldots,B_{\pi(r)}
\]
and unique delimiters $d_1,\ldots,d_{2r-1}$.  Define
\[
  P_{r,\pi}:=c_1^*d_1c_2^*d_2\cdots d_{2r-1}c_{2r}^*,
\]
and let $\delta$ map
\[
  A_i\mapsto e_i,
  \qquad
  B_i\mapsto -e_i,
  \qquad
  d_j\mapsto0
\]
in $\mathbb Z^r$.  Put
\[
  T_{r,\pi}:=P_{r,\pi}\cap\ker\delta.
\]

\begin{theorem}[Tagged Rambow--Satta witnesses are multidimensionally substitutable]
\label{thm:tagged-rs-s2}
For every $p\ge1$,
\[
  T_{r,\pi}\in S(p).
\]
In particular $T_{r,\pi}\in S(2)$.
\end{theorem}

\begin{proof}
Lemma~\ref{lem:unique-delimiter-chain} gives $P_{r,\pi}\in S(p)$, and
Lemma~\ref{lem:abelian-fiber-preservation} applied to the zero fiber of
$\delta:\Sigma^*\to\mathbb Z^r$ gives the result.
\end{proof}

\paragraph{Rank-only Rambow--Satta calibration.}
We make the parameter transfer used here explicit.  Consider the two-track
instance of the hierarchy witness in Rambow and Satta.  Number the blocks on
the first track from left to right by $1,\ldots,r$.  For every $r\ge6$, let
$\pi_r\in S_r$ be exactly the permutation fixed in
\cite[Definition~6]{RambowSatta1994}; Figure~3 of that report gives, for
orientation,
\[
  \pi_6(1),\ldots,\pi_6(6)=1,3,5,2,4,6,
  \qquad
  \pi_7(1),\ldots,\pi_7(7)=1,5,3,6,2,4,7.
\]
Thus $\pi_r$ is not an arbitrary permutation chosen later in the proof, but
the explicit hierarchy family of Rambow and Satta.  After renaming the two
track alphabets to $A_1,\ldots,A_r$ and $B_1,\ldots,B_r$, respectively,
their witness is
\[
 R^+_{r,\pi_r}=\{A_1^{n_1}\cdots A_r^{n_r}
 B_{\pi_r(1)}^{n_{\pi_r(1)}}\cdots B_{\pi_r(r)}^{n_{\pi_r(r)}}:
 n_1,\ldots,n_r\ge1\}.
\]
For every $r\ge6$, Rambow and Satta's hierarchy theorem gives the
independent-parallelism separation at levels $r-2$ and $r-3$
\cite[Theorem~1]{RambowSatta1994} (see also the refereed journal version
\cite{RambowSatta1999}).  Their language-class correspondence identifies the LUSCG and MCFG families
with the same fan-out and rank bounds \cite[Theorem~6]{RambowSatta1994}
(see also \cite{RambowSatta1999}).\footnote{Rambow and Satta write $q$-MCFG$(f)$ for
MCFGs with rule-rank bound $q$ and fan-out bound $f$; this is the language
family denoted $L(f,q)$ here.  On the LUSCG side, their rank is defined via
the right-hand-side length in the derivation grammar $\operatorname{der}(G)$
\cite[Definition~4]{RambowSatta1994}, rather than directly as the number of
nonterminals in an original sentential tuple.  The identification with MCFG
rule rank used here is supplied by their Theorem~6.}  Thus,
in the present notation,
\[
 R^+_{r,\pi_r}\in L(2,r-2)\setminus L(2,r-3).
\]
Fix this $\pi_r$ for each $r$, and define $R^0_{r,\pi_r}$ by allowing all
exponents to lie in $\mathbb N_0$.

For the closure steps below, Rambow and Satta prove that for every fixed rank
bound $q\ge2$ and fan-out bound $f\ge1$ the corresponding fixed-parameter
LUSCG family is a substitution-closed full AFL
\cite[Theorem~5]{RambowSatta1994} (see also \cite{RambowSatta1999}).
Theorem~6 of the technical report identifies it with the MCFG family having
those same two parameter bounds.  Consequently, in the parameter
range used below, $L(f,q)$ is closed \emph{within the same parameter stratum}
under homomorphism, inverse homomorphism, intersection with regular languages,
finite union, concatenation, Kleene plus, and substitution.  These are exactly
the closure properties used in the two transfers below.

\begin{proposition}[Exact minimum rank for the zero-exponent witness]
\label{prop:rs-zero-rank-only}
For every $r\ge6$,
\[
 R^0_{r,\pi_r}\in L(2,r-2)\setminus L(2,r-3).
\]
\end{proposition}

\begin{proof}
Let
\[
 \mathcal R_r=A_1^+\cdots A_r^+
 B_{\pi_r(1)}^+\cdots B_{\pi_r(r)}^+.
\]
Then $R^+_{r,\pi_r}=R^0_{r,\pi_r}\cap\mathcal R_r$.  If
$R^0_{r,\pi_r}$ had a fan-out-two presentation of rule rank at most $r-3$,
closure under regular intersection within the same fixed fan-out/rule-rank
stratum would contradict the displayed Rambow--Satta separation.

For the converse, for each $Z\subseteq\{1,\ldots,r\}$ let $e_Z$ erase both
$A_i$ and $B_i$ for $i\in Z$ and fix all remaining terminals.  Then
\[
 R^0_{r,\pi_r}=\bigcup_{Z\subseteq\{1,\ldots,r\}} e_Z(R^+_{r,\pi_r}).
\]
The fixed-parameter homomorphism and finite-union closures stated above
therefore transfer the upper bound $r-2$ from $R^+_{r,\pi_r}$ without
changing fan-out or rule rank.
\end{proof}

Let $e$ erase all delimiters $d_j$ and fix the phase letters.  For the tagged
language from Section~\ref{sec:comparison},
\[
 e(T_{r,\pi})=R^0_{r,\pi},
 \qquad
 T_{r,\pi}=e^{-1}(R^0_{r,\pi})\cap P_{r,\pi}.
\]
By Rambow and Satta's fixed-parameter full-AFL result and their
LUSCG--MCFG correspondence \cite[Theorems~5 and~6]{RambowSatta1994}
(see also \cite{RambowSatta1999}), inverse homomorphism by $e$ and intersection
with the regular language $P_{r,\pi}$ preserve the same fan-out and
rule-rank stratum in the parameter range at hand.  Therefore
Proposition~\ref{prop:rs-zero-rank-only}
transfers the exact rank $r-2$ to $T_{r,\pi_r}$.  On the other hand,
Theorem~\ref{thm:tagged-rs-s2} gives $T_{r,\pi_r}\in S(2)$.  These two facts
complete the proof of Theorem~\ref{thm:s2-rank-hierarchy}.


\section{Decision Complexity of the Semantic Promise}
\label{app:decision-complexity}

The learning algorithms do not need to recognize the semantic promise from a
target grammar.  This matters also for the use of finite characteristic
samples: the natural grammar presentation cannot simply be filtered
effectively into an indexed family for the semantic slice.  For completeness,
this appendix records the exact computability boundary of that recognition
problem and the contrasting regular case.


The learner never receives a target grammar, so the learnability results above
do not require an algorithm for recognizing membership in a semantic slice.
Nevertheless, the recognition problem itself gives a useful decision-theoretic
boundary for the framework.  We use $\Pi^0_1$ synonymously with
co-recursively enumerable (co-r.e.).

\begin{proposition}[The promise-recognition problem is co-r.e.]
\label{prop:promise-pi1-upper}
Fix $f\ge1$.  Given a standard $f$-MCFG $G$ over $\Sigma$ and an explicit
finite observation $h:\Sigma^*\to M$, nonmembership of $L(G)$ in
$\Cmcf{f}{h}$ is recursively enumerable.  Equivalently, recognition of
$(f,h)$-tuple-substitutability belongs to $\Pi^0_1$.
\end{proposition}

\begin{proof}
A failure of Definition~\ref{def:tuple-substitutable} has a finite witness:
some $d\le f$, orientation $\sigma\in S_d$, tuples
$\tuple x,\tuple y\in(\Sigma^*)^d$, and named sentence contexts
$E,F\in\mathcal E_{d,\sigma}$ such that
\[
 h^{(d)}(\tuple x)=h^{(d)}(\tuple y),
 \qquad
 E[\tuple x],E[\tuple y]\in L(G),
\]
and, after exchanging $\tuple x,\tuple y$ if necessary,
\[
 F[\tuple x]\in L(G),
 \qquad
 F[\tuple y]\notin L(G).
\]
All of these finite objects can be effectively enumerated.  Membership in an
MCFG language is decidable, and the explicit finite observation $h$ is
effectively evaluable.  Hence failures can be enumerated.
\end{proof}

\begin{theorem}[Semantic-promise recognition is $\Pi^0_1$-complete]
\label{thm:promise-pi1-complete}
For every fixed $f\ge1$, the following problem is $\Pi^0_1$-complete:
\begin{quote}
Given a standard $f$-MCFG $G$ and an explicit finite observation $h$, decide
whether $L(G)$ is $(f,h)$-tuple-substitutable.
\end{quote}
Hardness already holds when $G$ is restricted to context-free grammars and
$h=\mathbf 1$ is the one-element observation.
\end{theorem}

\begin{proof}
The upper bound is Proposition~\ref{prop:promise-pi1-upper}.  For hardness we
use context-free universality.  The classical computation-history construction
is a many-one reduction from nonacceptance: from a Turing machine $T$ and input
$x$ it effectively constructs a CFG $G_{T,x}$ such that
\[
  L(G_{T,x})=\Delta^*
  \quad\Longleftrightarrow\quad
  T\text{ does not accept }x
\]
\cite[p.~281]{HopcroftUllman1979}.  Hence
$\overline{A_{\rm TM}}\le_m\mathrm{UNIV}_{\rm CFG}$, so CFG universality is
$\Pi^0_1$-hard.  Conversely, nonuniversality is recursively enumerable:
enumerate the words of $\Delta^*$ and use decidable CFG membership until a
missing word is found.  Thus universality is co-r.e., equivalently in
$\Pi^0_1$, and is therefore $\Pi^0_1$-complete.  Let $G_0$ be a CFG over
an alphabet $\Delta$ with
$|\Delta|\ge2$, write $L_0=L(G_0)$, choose a fresh letter
$d\notin\Delta$, and put $\Sigma=\Delta\uplus\{d\}$.

First test the decidable conditions $\eps\in L_0$ and
$L_0\cap\Delta^+\ne\varnothing$.  If either fails, then $L_0\ne\Delta^*$ is
already known; output a fixed CFG for the non-substitutable language
$\{a,aa\}$, where $a\in\Delta$.  Otherwise construct effectively a CFG for
\[
 L' := L_0\ \cup\ D,
 \qquad
 D:=\Sigma^*d\Sigma^*.
\]
If $L_0=\Delta^*$, every word over $\Sigma$ either contains $d$ or belongs to
$\Delta^*$, so $L'=\Sigma^*$.  Hence $L'$ is $(f,\mathbf1)$-tuple-substitutable
for every $f$.

Suppose instead that $L_0\ne\Delta^*$.  Under the branch assumptions there
exist
\[
 y\in L_0\cap\Delta^+,
 \qquad
 z\in\Delta^+\setminus L_0.
\]
At arity one the strings $y,z$ have the same trivial observation type.  They
share the accepting context $E=d\blank_1$, since both $dy$ and $dz$ lie in
$D\subseteq L'$.  Their distributions differ in the empty context
$F=\blank_1$: one has $F[y]=y\in L'$, whereas $F[z]=z\notin L'$ because $z$
contains no $d$ and $z\notin L_0$.  Thus $L'$ is not
$(f,\mathbf1)$-tuple-substitutable.  Consequently
\[
 L'\text{ is }(f,\mathbf1)\text{-tuple-substitutable}
 \quad\Longleftrightarrow\quad
 L_0=\Delta^*,
\]
which proves $\Pi^0_1$-hardness.  The fixed fallback language $\{a,aa\}$ is
indeed non-substitutable: $a$ and $aa$ share the empty context, while
$a\blank_1$ accepts $a$ and rejects $aa$.
\end{proof}


\begin{thebibliography}{99}

\bibitem{Angluin1980}
D.~Angluin.
\newblock Inductive inference of formal languages from positive data.
\newblock {\em Information and Control}, 45(2):117--135, 1980.

\bibitem{Angluin1982Reversible}
D.~Angluin.
\newblock Inference of reversible languages.
\newblock {\em Journal of the ACM}, 29(3):741--765, 1982.

\bibitem{CaseYoshinakaZeugmann2013}
J.~Case, R.~Yoshinaka, and T.~Zeugmann.
\newblock Stochastic finite learning of some mildly context-sensitive languages.
\newblock Presented at the ICALP 2013 Satellite Workshop on Learning Theory and
Complexity, Riga, Latvia, 2013.
\newblock Author manuscript available at
\url{https://www-alg.ist.hokudai.ac.jp/~thomas/LTC/PAPERS/yoshinaka.pdf}.

\bibitem{CaseLynes1982}
J.~Case and C.~Lynes.
\newblock Machine inductive inference and language identification.
\newblock In M.~Nielsen and E.~M. Schmidt (eds.), \emph{Automata, Languages
and Programming: 9th Colloquium}, Lecture Notes in Computer Science, vol.~140,
pp.~107--115. Springer, 1982.
\newblock doi:10.1007/BFb0012761.

\bibitem{CaseSmith1983}
J. Case and C. Smith,
``Comparison of identification criteria for machine inductive inference,''
\emph{Theoretical Computer Science},
Vol.~25, No.~2, 1983, pp.~193--220.
doi:10.1016/0304-3975(83)90061-0.

\bibitem{Clark2017}
A.~Clark.
\newblock Testing distributional properties of context-free grammars.
\newblock In {\em Proceedings of the 13th International Conference on
Grammatical Inference}, PMLR 57, pp.~42--53, 2017.

\bibitem{Clark2021Strong}
A.~Clark.
\newblock Strong learning of some probabilistic multiple context-free grammars.
\newblock In {\em Proceedings of the 17th Meeting on the Mathematics of Language},
pp.~23--37. Association for Computational Linguistics, 2021.

\bibitem{ClarkEyraud2007}
A.~Clark and R.~Eyraud.
\newblock Polynomial identification in the limit of substitutable context-free
languages.
\newblock {\em Journal of Machine Learning Research}, 8:1725--1745, 2007.

\bibitem{ClarkYoshinaka2012PMCFG}
A.~Clark and R.~Yoshinaka.
\newblock Beyond semilinearity: Distributional learning of parallel multiple
context-free grammars.
\newblock In {\em Proceedings of the Eleventh International Conference on
Grammatical Inference}, PMLR 21, pp.~84--96, 2012.

\bibitem{ClarkYoshinaka2014}
A.~Clark and R.~Yoshinaka.
\newblock Distributional learning of parallel multiple context-free grammars.
\newblock {\em Machine Learning}, 96(1--2):5--31, 2014.
\newblock doi:10.1007/s10994-013-5403-2.


\bibitem{ClarkYoshinaka2014Algebraic}
A.~Clark and R.~Yoshinaka.
\newblock An algebraic approach to multiple context-free grammars.
\newblock In {\em Logical Aspects of Computational Linguistics: 8th International
Conference, LACL 2014}, LNCS 8535, pp.~57--69. Springer, 2014.
\newblock doi:10.1007/978-3-662-43742-1\_5.

\bibitem{ClarkYoshinaka2016}
A.~Clark and R.~Yoshinaka.
\newblock Distributional learning of context-free and multiple context-free
grammars.
\newblock In {\em Topics in Grammatical Inference}, pp.~143--172. Springer,
2016.
\newblock doi:10.1007/978-3-662-48395-4\_6.

\bibitem{Gold1967}
E.~M.~Gold.
\newblock Language identification in the limit.
\newblock {\em Information and Control}, 10(5):447--474, 1967.

\bibitem{GomezRodriguezEtAl2009}
C.~G\'omez-Rodr\'iguez, M.~Kuhlmann, G.~Satta, and D.~Weir.
\newblock Optimal reduction of rule length in linear context-free rewriting systems.
\newblock In {\em Proceedings of Human Language Technologies: The 2009 Annual
Conference of the North American Chapter of the Association for Computational
Linguistics}, pp.~539--547. Association for Computational Linguistics, 2009.

\bibitem{HopcroftUllman1979}
J.~E. Hopcroft and J.~D. Ullman.
\newblock \emph{Introduction to Automata Theory, Languages, and Computation}.
\newblock Addison--Wesley, 1979.

\bibitem{Kallmeyer2010}
L.~Kallmeyer.
\newblock {\em Parsing Beyond Context-Free Grammars}.
\newblock Springer, Berlin, Heidelberg, 2010.
\newblock doi:10.1007/978-3-642-14846-0.

\bibitem{Kanazawa2019Ogden}
M.~Kanazawa.
\newblock Ogden's lemma, multiple context-free grammars, and the control language hierarchy.
\newblock {\em Information and Computation}, 269:104449, 2019.
\newblock doi:10.1016/j.ic.2019.104449.

\bibitem{KanazawaKappe2019}
M.~Kanazawa and T.~Kapp\'{e}.
\newblock Decision problems for Clark-congruential languages.
\newblock In {\em Proceedings of the 14th International Conference on
Grammatical Inference}, PMLR 93, pp.~3--16, 2019.

\bibitem{KanazawaYoshinaka2017}
M.~Kanazawa and R.~Yoshinaka.
\newblock The strong, weak, and very weak finite context and kernel properties.
\newblock In {\em Language and Automata Theory and Applications}, LNCS 10168,
pp.~77--88. Springer, 2017.
\newblock doi:10.1007/978-3-319-53733-7\_5.

\bibitem{KanazawaYoshinaka2021}
M.~Kanazawa and R.~Yoshinaka.
\newblock A hierarchy of context-free languages learnable from positive data
and membership queries.
\newblock In {\em Proceedings of the 15th International Conference on
Grammatical Inference}, PMLR 153, pp.~18--31, 2021.

\bibitem{KanazawaYoshinaka2023}
M.~Kanazawa and R.~Yoshinaka.
\newblock Extending distributional learning from positive data and membership
queries.
\newblock In {\em Proceedings of the 16th International Conference on
Grammatical Inference}, PMLR 217, pp.~8--22, 2023.

\bibitem{kuriyamaCFG}
T.~Kuriyama.
\newblock Distributional learning of context-free languages under fixed finite-monoid typing.
\newblock arXiv:1409.6247 [cs.FL], 2014; revised version v4, May 2026.

\bibitem{Pin1986}
J.-E. Pin.
\newblock {\em Varieties of Formal Languages}.
\newblock North Oxford Academic, London, and Plenum, New York, 1986.

\bibitem{RambowSatta1994}
O.~Rambow and G.~Satta.
\newblock A two-dimensional hierarchy for parallel rewriting systems.
\newblock Technical Report IRCS-94-02, Institute for Research in Cognitive
Science, University of Pennsylvania, 1994.
\newblock Public repository copy:
\url{https://repository.upenn.edu/ircs_reports/148/}.

\bibitem{RambowSatta1999}
O.~Rambow and G.~Satta.
\newblock Independent parallelism in finite copying parallel rewriting systems.
\newblock {\em Theoretical Computer Science}, 223(1--2):87--120, 1999.
\newblock doi:10.1016/S0304-3975(97)00190-4.

\bibitem{SagotSatta2010}
B.~Sagot and G.~Satta.
\newblock Optimal rank reduction for linear context-free rewriting systems with fan-out two.
\newblock In {\em Proceedings of the 48th Annual Meeting of the Association for
Computational Linguistics}, pp.~525--533. Association for Computational Linguistics, 2010.

\bibitem{SekiEtAl1991}
H.~Seki, T.~Matsumura, M.~Fujii, and T.~Kasami.
\newblock On multiple context-free grammars.
\newblock {\em Theoretical Computer Science}, 88(2):191--229, 1991.
\newblock doi:10.1016/0304-3975(91)90374-B.

\bibitem{Yoshinaka2008}
R.~Yoshinaka.
\newblock Identification in the limit of \(k,\ell\)-substitutable context-free
 languages.
\newblock In {\em Grammatical Inference: Algorithms and Applications},
 LNCS 5278, pp.~266--279. Springer, 2008.
\newblock doi:10.1007/978-3-540-88009-7\_21.

\bibitem{Yoshinaka2009}
R.~Yoshinaka.
\newblock Learning mildly context-sensitive languages with multidimensional
 substitutability from positive data.
\newblock In {\em Algorithmic Learning Theory}, LNCS 5809, pp.~278--292.
 Springer, 2009.
\newblock doi:10.1007/978-3-642-04414-4\_24.

\bibitem{Yoshinaka2010}
R.~Yoshinaka.
\newblock Polynomial-time identification of multiple context-free languages
from positive data and membership queries.
\newblock In {\em Grammatical Inference: Theoretical Results and Applications},
LNCS 6339, pp.~230--244. Springer, 2010.
\newblock doi:10.1007/978-3-642-15488-1\_19.

\bibitem{Yoshinaka2011}
R.~Yoshinaka.
\newblock Efficient learning of multiple context-free languages with
multidimensional substitutability from positive data.
\newblock {\em Theoretical Computer Science}, 412(19):1821--1831, 2011.
\newblock doi:10.1016/j.tcs.2010.12.058.

\bibitem{YoshinakaClark2012}
R.~Yoshinaka and A.~Clark.
\newblock Polynomial time learning of some multiple context-free languages
with a minimally adequate teacher.
\newblock In {\em Formal Grammar}, LNCS 7395, pp.~192--207. Springer, 2012.
\newblock doi:10.1007/978-3-642-32024-8\_13.

\end{thebibliography}

\end{document}
